% !TeX TXS-program:bibliography = txs:///biber
% !BIB program = biber
\documentclass[11pt,letterpaper]{article}

\usepackage[T1]{fontenc}
\usepackage{lmodern}
\usepackage{microtype}
\usepackage{amsmath,amssymb,amsthm,mathtools}
\usepackage[margin=1.03in]{geometry}
\usepackage{xcolor}
\usepackage{tikz}
\definecolor{DarkTurquoise}{HTML}{006D77}
\definecolor{ReferenceViolet}{HTML}{6F2DA8}
\usepackage[
  backend=biber,
  style=alphabetic,
  sorting=nyt,
  giveninits=true,
  maxbibnames=99,
  doi=false,
  isbn=false,
  url=false,
  eprint=false
]{biblatex}
\addbibresource{cases_234_codex.bib}
\AtBeginBibliography{\setlength{\emergencystretch}{2em}}
\usepackage{xurl}
\usepackage{hyperref}
\hypersetup{
  colorlinks=true,
  linkcolor=DarkTurquoise,
  citecolor=ReferenceViolet,
  urlcolor=DarkTurquoise,
  pdfauthor={Pablo Portilla Cuadrado},
  pdftitle={Injective holomorphic map germs of low order}
}
\usepackage{aliascnt}
\usepackage[capitalise,nameinlink,noabbrev]{cleveref}

\numberwithin{equation}{section}
\newaliascnt{theorem}{equation}
\newtheorem{theorem}[theorem]{Theorem}
\aliascntresetthe{theorem}
\newaliascnt{proposition}{equation}
\newtheorem{proposition}[proposition]{Proposition}
\aliascntresetthe{proposition}
\newaliascnt{lemma}{equation}
\newtheorem{lemma}[lemma]{Lemma}
\aliascntresetthe{lemma}
\newaliascnt{corollary}{equation}
\newtheorem{corollary}[corollary]{Corollary}
\aliascntresetthe{corollary}
\theoremstyle{remark}
\newaliascnt{remark}{equation}
\newtheorem{remark}[remark]{Remark}
\aliascntresetthe{remark}
\theoremstyle{plain}

\newcommand{\C}{\mathbb{C}}
\newcommand{\Z}{\mathbb{Z}}
\newcommand{\ord}{\operatorname{ord}}
\newcommand{\mult}{\operatorname{mult}}

\title{Injective analytic mappings of low order}
\author{Pablo Portilla Cuadrado}
\date{\today}

\AtBeginDocument{
  \setlength{\abovedisplayskip}{9pt plus 3pt minus 6pt}
  \setlength{\belowdisplayskip}{9pt plus 3pt minus 6pt}
}

\begin{document}
\maketitle

\begin{abstract}
Let $\boldsymbol n:(\C^2,0)\to(\C^3,0)$ be a holomorphic map germ
admitting an injective representative.  We prove that
$(d\boldsymbol n)_0\neq0$, as conjectured by L\^e.
Building on the author's earlier reduction theorem
\cite{PortillaReduction}, this work completes the proof of the conjecture by
excluding maps of orders $2$, $3$, and $4$.
The proofs combine successive analytic changes of source and target
coordinates, Weierstrass preparation, and the analysis of weighted
initial forms with techniques from singularity theory, including
semicontinuity of the spectrum. In dialogue with the author, ChatGPT
developed most of the mathematical ideas, repaired a gap in the author's
initial argument for order $4$, and wrote most of the manuscript subsequently 
refined by the author.
\end{abstract}

\section{Introduction}

Let
\[
 \boldsymbol n:(\C^2,0)\longrightarrow(\C^3,0),\qquad
 (x,y)\longmapsto(n_1(x,y),n_2(x,y),n_3(x,y))
\]
be a holomorphic map germ admitting an injective representative.
Its \emph{order} is
$
 m_{\boldsymbol n}=\ord_0(\boldsymbol n)
 =\min_{1\leq j\leq3}\ord_0(n_j).
$
The reduction theorem \cite[Theorem~1.1]{PortillaReduction} states that
$(d\boldsymbol n)_0=0$ implies
$m_{\boldsymbol n}\in\{2,3,4\}$.  The purpose of this paper is to
exclude these three orders and thereby prove L\^e's conjecture.
We refer to the introduction of
\cite{PortillaReduction} for the geometric formulation in terms of
links of reduced surface germs, the equivalence of the two formulations,
and the history of the problem.

\begin{theorem}
\label{thm:main}
Every holomorphic map germ
$
 \boldsymbol n:(\C^2,0)\longrightarrow(\C^3,0)
$
that admits an injective representative satisfies
$(d\boldsymbol n)_0\neq0$.
\end{theorem}

\subsection*{Strategy and organization of the proof}
\label{sec:strategy}

Suppose that $(d\boldsymbol n)_0=0$.  Choose linear coordinates on
the target as in \Cref{sec:input}, and write
$\boldsymbol n=(v,u,w)$ in these coordinates, so that $(v,u)$ is finite.
Put $X=\boldsymbol n(\C^2)$, $V=\{v=0\}$,
$m_{\boldsymbol n}\in\{2,3,4\}$, and $m_X=\mult_0(X)$.
The geometric question is whether the third coordinate can distinguish
all the points with the same first two coordinates.

To see this on a curve, choose $\delta>0$ sufficiently small and then
$\eta>0$, and put $\overline D_\delta=\{z\in\C:|z|\leq\delta\}$.
For $0<|s|<\eta$, choose a Milnor representative
$V_s$ of $\{v=s\}$ so that
\[
 u:V_s\longrightarrow\overline D_\delta,\qquad p\longmapsto u(p)
\]
is a proper branched covering of degree $m_X$, unramified on the
boundary, as in \cite[proof of Proposition~2.16]{PortillaReduction}.
This is the covering used to obtain the upper bound on transverse
multiplicity in the reduction theorem.  Write $W_s=\boldsymbol n(V_s)$
for the corresponding representative of the target section.  On $V_s$,
the map $\boldsymbol n$ is
\[
 \boldsymbol n|_{V_s}:V_s\longrightarrow W_s,\qquad
 p\longmapsto(s,u(p),w(p)).
\]
For $p,q\in V_s$, the first coordinates agree.  Consequently,
\begin{equation}
 \boldsymbol n(p)=\boldsymbol n(q)
 \quad\Longleftrightarrow\quad
 u(p)=u(q)\ \text{and}\ w(p)=w(q).
 \label{eq:strategy-separation}
\end{equation}
Thus injectivity requires $w$ to take distinct values at every pair of
distinct points in a fiber of $u|_{V_s}$.
We call a pair of distinct points with the same image under a map
a \emph{collision} of that map.
The same requirement applies to every fiber of the finite map $(v,u)$.

The proof develops this separation requirement in three steps.
\begin{enumerate}
\item \label{step:ramification}
\emph{Count the ramification and prepare a critical component.}
Writing $\mu(V,0)$ for the Milnor number of $V$, the Riemann--Hurwitz
formula gives total ramification $m_X+\mu(V,0)-1$ for $u|_{V_s}$.
Choose a component $\Gamma_i$ of the critical locus whose image has
maximal transverse multiplicity $m_{\mathrm{tr}}$, and put
$a_i=i_0(v,\Gamma_i)$, as in \Cref{sec:input}.
The numerical inequalities and base-point formulas yield
\Cref{prop:common-preparation}: $\Gamma_i$ is smooth and
$a_i=m_{\boldsymbol n}$.
The component $\Gamma_i$ meets $V_s$ in $a_i$ points.  At each of them
the image has a transverse branch of multiplicity $m_{\mathrm{tr}}$,
so the local degree of $u$ is at least $m_{\mathrm{tr}}$.
These points contribute at least $a_i(m_{\mathrm{tr}}-1)$ to the
total ramification, giving
\[
 a_i(m_{\mathrm{tr}}-1)\leq m_X+\mu(V,0)-1.
\]
The same transverse multiplicity forces the Jacobian minors of
$\boldsymbol n$ to vanish to order at least $m_{\mathrm{tr}}-1$
along $\Gamma_i$.

\item \label{step:root-separation}
\emph{Determine the first separation of the roots.}
Make $\Gamma_i$ the axis $y=0$ and write $V$ as the zero set of a
Weierstrass polynomial of degree $m_{\boldsymbol n}$ in $x$, with the
roots translated so that their mean is zero.  Rescale at the first order
at which its roots separate as $y\to0$.  The resulting weighted
initial forms must satisfy both the numerical inequalities and
the divisibility of the Jacobian minors from
Step~\ref{step:ramification}.  These restrictions exclude order four
in \Cref{sec:dfour}.  In order three they leave two possible pairs
of initial forms.  Order two requires a further analytic preparation.
Here $m_{\boldsymbol n}$ counts the roots of the Weierstrass equation.  The covering
$u|_{V_s}$ has degree $m_X$, which can be much larger.

\item \label{step:third-coordinate}
\emph{Show that the third coordinate fails to separate a pair.}
For the remaining orders, study pairs of points with the same first
two coordinates and the difference of their third coordinates.
Subtracting a function of the first two coordinates does not change
this difference.  When the first two initial forms define a finite
map, \Cref{lem:filtered-escape} uses such subtractions to expose a
third initial form outside their polynomial algebra.
The fiber products of the pairs of initial forms describe the pairs
of points that the third initial form must separate.  In order three,
symmetries and a derivative identity force equal third-coordinate values at distinct
points.  In the other
projective curve arguments, the difference is a section whose zeros
are counted with multiplicity, and the bounds prevent all its zeros
from lying on the diagonal.  An isolated zero corresponding to
distinct points lifts, by \Cref{lem:weighted-lifting}, to collisions
of the analytic map arbitrarily close to the origin.
Order two also uses a local fiber product disjoint from the diagonal.
Injectivity would confine the zeros of the third-coordinate difference
to its limiting fiber.  A second rescaling shows that this constraint
forces a collision.
Each collision violates \eqref{eq:strategy-separation}.
\end{enumerate}

\Cref{sec:auxiliary} establishes the common preparation, the weighted
estimates, and the lifting criterion used in these steps.  The three
subsequent sections carry out these steps in orders four, three, and two,
respectively.

\subsection*{AI disclosure}

This paper was essentially written by ChatGPT,
through several conversations and extended, largely autonomous work in
dialogue with the author. ChatGPT developed the proofs for
$m_{\boldsymbol n}=2,3,4$, including the repair of a gap in the author's
initial argument for $m_{\boldsymbol n}=4$. The author's contribution to this 
paper has consisted primarily of
refining and substantially shortening the generated text, organizing
the ideas, and reviewing and verifying the arguments. The author takes
responsibility for all the arguments in the final manuscript.

The proofs below rely, sometimes, on the results of the companion manuscript
\emph{A reduction theorem for L\^e's conjecture}
\cite{PortillaReduction} which was part of the prompt that ChatGPT used to 
solve these cases. In the course of this work, we have not yet
found a proof of the cases $m_{\boldsymbol n}=2,3,4$ that does not rely
on that work.

Although the proofs of the three cases are presented in a common
framework, some auxiliary results remain specific to individual cases.
Examples include \Cref{prop:d3-cubic-collision,prop:d2-cyclic-separation}.
The proof of \Cref{prop:d2-cyclic-separation} is long because its steps
depend closely on the particular form of the map. This dependence makes
it difficult to divide the proof into shorter arguments that can be
reused elsewhere in the paper. There are others like this.

We draw special attention to \Cref{lem:d2-forbidden-splitting}, whose proof
uses semicontinuity of the spectrum and whose statement will trap a specialist 
in plane curve singularities an afternoon or two. We have not found an 
elementary
proof of this lemma.

\section{Input from the reduction theorem}
\label{sec:input}


We use the notation of the reduction theorem.  Throughout the proof, let
$\boldsymbol n=(n_1,n_2,n_3):(\C^2,0)\to(\C^3,0)$ be a holomorphic
map germ admitting an injective representative and satisfying
$(d\boldsymbol n)_0=0$.  Write $(X,0)$ for its reduced image germ.
We fix an injective representative of $\boldsymbol n$ and a compatible
representative of $X$, denoted by the same symbols.  Choose coordinates
$(x,y)$ on $(\C^2,0)$ and put
\[
 B=\C\{x,y\},\qquad \mathfrak m=(x,y)B,\qquad
 I=(n_1,n_2,n_3)B,\qquad m_{\boldsymbol n}=\ord_0 I.
\]
The map $\boldsymbol n:(\C^2,0)\to(X,0)$ is the normalization map, and
$I$ is $\mathfrak m$-primary.  Starting with $I$, at each blowup we
pull back local generators of the current ideal and divide them by
their highest common power of a local equation of the new exceptional
curve.  The ideal generated by these quotients is the new
\emph{weak transform}.  We blow up the points where this ideal is not
the unit ideal, and stop when it is the unit ideal.  This is the
\emph{principalization} used in the reduction theorem.
Let $\mathcal C$ be the set of
centres of these blowups, called the \emph{cluster of base points}, and
let $m_{\boldsymbol n,p}$ be the order of the weak transform at $p$,
as in \cite[Notation~2.4]{PortillaReduction}.
For $f,g\in B$, write $J(f,g)=f_xg_y-f_yg_x$ for their
\emph{Jacobian determinant}
with respect to $(x,y)$.  For other coordinates $(s,t)$, we write
$J_{s,t}(f,g)=f_sg_t-f_tg_s$.
The subscript records the differentiation variables in order, so
$J_{t,s}(f,g)=-J_{s,t}(f,g)$.
For $f,g\in B$ whose zero sets have no common curve component, their
\emph{intersection multiplicity} at the origin is
\[
 i_0(f,g)=\dim_\C B/(f,g)B.
\]
For $f\in B$ and a reduced plane curve germ $\Gamma\subset(\C^2,0)$
having no component in $\{f=0\}$, choose a reduced equation $h=0$ for $\Gamma$
and set $i_0(f,\Gamma)=i_0(f,h)$.

Choose general linear forms $\lambda_1,\lambda_2$ on the target as in
\cite[Sections~2.1--2.2]{PortillaReduction}, and complete them to target
coordinates with a third linear form $\lambda_3$.  Put
$v=\lambda_1\circ\boldsymbol n$, $u=\lambda_2\circ\boldsymbol n$,
and $w=\lambda_3\circ\boldsymbol n$.
Thus $V=\{v=0\}$
is reduced, its strict transform has multiplicity
$m_{\boldsymbol n,p}$ at each $p\in\mathcal C$, and its final
strict transform is smooth, transverse to the exceptional divisor,
and disjoint from its crossings.  There are no infinitely near
singular points of $V$ outside $\mathcal C$.  Moreover,
$\{v=u=0\}=\{0\}$, and $u$ has minimal order on each branch
of $V$ among the elements of $I$.
Write $W=X\cap\{\lambda_1=0\}$, and denote the \emph{delta invariant}
and \emph{Milnor number} of a reduced plane curve germ by $\delta$ and
$\mu$, respectively, following
\cite[p.~206 and Definition~I.2.1]{GreuelLossenShustin}.
Let $b$ be the number of branches of $V$, let
$m_X=\mult_0(X)=\mult_0(W)$, and put
$M_{\boldsymbol n}=m_X-2\delta(V,0)$.

Write $\Sigma=\operatorname{Sing}(X)=\bigcup_{i=1}^c\Sigma_i$
for the decomposition into irreducible components, and put
$\Gamma_i=(\boldsymbol n^{-1}(\Sigma_i))_{\mathrm{red}}$.
The curves $\Gamma_i$ are the irreducible components of the critical
locus of $\boldsymbol n$.  As in the reduction theorem, the
\emph{transverse multiplicity} $m_i$ along $\Sigma_i$ is the multiplicity
of a transverse plane branch of $X$ at a general point of $\Sigma_i$.  Set
\[
 a_i=i_0(v,\Gamma_i),\qquad m_{\mathrm{tr}}=\max_i m_i\geq2.
\]
We also choose $v$ so that the final strict transform of $V$ avoids the
strict transforms of the finitely many $\Gamma_i$ over the origin.
The strict transforms of the $\Gamma_i$ meet the exceptional divisor
in finitely many points.  After dividing by their common exceptional
factors, the pullbacks of $n_1,n_2,n_3$ have no common zero.  Thus
requiring the strict transform of $V$ to pass through any one of these
points imposes a nontrivial linear equation on the coefficients of
$\lambda_1$.
A general choice of $v$ avoids all these conditions. We use the following 
identities from
\cite[equations~(2.18)--(2.19)]{PortillaReduction}:
\begin{equation}
 m_X=\sum_{p\in\mathcal C}m_{\boldsymbol n,p}^2,\qquad
 M_{\boldsymbol n}=\sum_{p\in\mathcal C}m_{\boldsymbol n,p}.
 \label{eq:input-point-basis}
\end{equation}
The same choice of $v,u$ gives
\begin{equation}
 i_0(v,u)=m_X,\qquad
 2\delta(V,0)=m_X-M_{\boldsymbol n}.                 \label{eq:input-point}
\end{equation}
The reduction inequalities
\cite[Propositions~2.16 and~2.17, Theorem~4.5]{PortillaReduction} are
\begin{align}
 M_{\boldsymbol n}&\geq m_{\boldsymbol n},&
 m_X&\leq m_{\boldsymbol n}M_{\boldsymbol n},          \label{eq:input-resolution}\\
 M_{\boldsymbol n}&\leq m_{\mathrm{tr}},&
 m_X+M_{\boldsymbol n}&\leq3m_{\mathrm{tr}},            \label{eq:input-knot}\\
 \sum_i a_i(m_i-1)&\leq2m_X-M_{\boldsymbol n}-b.        \label{eq:input-polar}
\end{align}
The plane curve identity $\mu=2\delta-b+1$
\cite[Proposition~I.3.35]{GreuelLossenShustin}, applied to $(V,0)$, gives
\begin{equation}
 m_X-M_{\boldsymbol n}=\mu(V,0)+b-1.                 \label{eq:input-milnor}
\end{equation}
The polar bound \eqref{eq:input-polar} is the ramification count in
Step~\ref{step:ramification} of the \hyperref[sec:strategy]{strategy}.
For the covering $u:V_s\to\overline D_\delta$ described there,
Riemann--Hurwitz and \eqref{eq:input-milnor} give total ramification
$m_X+\mu(V,0)-1=2m_X-M_{\boldsymbol n}-b$.
The $a_i$ points of $V_s\cap\Gamma_i$ each contribute at least
$m_i-1$.  Summing these contributions gives \eqref{eq:input-polar}.
The polar bound constrains the first two coordinates $v,u$.  By
\eqref{eq:strategy-separation}, the third coordinate $w$ must take
different values at distinct points with the same values of $v,u$.

Choose reduced equations $h_i=0$ for the curves $\Gamma_i$.  The
transverse multiplicities $m_i$ along $\Sigma_i$ give
\begin{equation}
 \prod_i h_i^{m_i-1}\mid J(n_j,n_k)
 \qquad(1\leq j<k\leq3).                               \label{eq:input-divisor}
\end{equation}
Indeed, at a general point of $\Gamma_i$, take source coordinates
$(s,z)$ and target coordinates in which
$\boldsymbol n=(s,a(s,z),c(s,z))$ and $\Gamma_i=\{z=0\}$.
The series $a(s,z)-a(s,0)$ and $c(s,z)-c(s,0)$ have
$z$-order at least $m_i$, so all three Jacobian minors of
$\boldsymbol n$ are divisible by $z^{m_i-1}$.  The orders of these
minors along the irreducible curves $\Gamma_i$, together with unique
factorization in $B$, prove \eqref{eq:input-divisor}.

Choose $i$ with $m_i=m_{\mathrm{tr}}$.  We next bound
$a_i=i_0(v,\Gamma_i)$.  This bound will imply that $\Gamma_i$ is smooth
when $\boldsymbol n$ has order three or four.  The local intersection
inequality and \eqref{eq:input-polar} give
\begin{equation}
 a_i\geq m_{\boldsymbol{n}},\qquad
 a_i(m_{\mathrm{tr}}-1)
 \leq2m_X-M_{\boldsymbol{n}}-b.                        \label{eq:one-polar}
\end{equation}

\begin{lemma}
\label{lem:maximal-component-estimate}
Suppose that $m_{\boldsymbol{n}}\in\{3,4\}$.  For every index $i$ with
$m_i=m_{\mathrm{tr}}$, one has $a_i\leq4$.  If $a_i=4$, then
\begin{equation}
 2m_X\geq7M_{\boldsymbol{n}}+3b-12.                    \label{eq:maximal-component-four}
\end{equation}
\end{lemma}

\begin{proof}
If $a_i\geq5$, then \eqref{eq:input-knot} and \eqref{eq:one-polar} give
\[
 5\left(\frac{m_X+M_{\boldsymbol{n}}}{3}-1\right)
 \leq2m_X-M_{\boldsymbol{n}}-b,
\]
and therefore, by \eqref{eq:input-resolution} and $b\geq1$,
\[
 \begin{aligned}
 0&\leq m_X-8M_{\boldsymbol n}+15-3b\\
  &\leq(m_{\boldsymbol n}-8)M_{\boldsymbol n}+12\\
  &\leq(m_{\boldsymbol n}-8)m_{\boldsymbol n}+12<0.
 \end{aligned}
\]
The last expression is $-3$ when $m_{\boldsymbol n}=3$ and $-4$
when $m_{\boldsymbol n}=4$.
Thus $a_i\leq4$. If $a_i=4$, \eqref{eq:input-knot} and
\eqref{eq:one-polar} give
\[
 4\left(\frac{m_X+M_{\boldsymbol{n}}}{3}-1\right)
 \leq4(m_{\mathrm{tr}}-1)
 \leq2m_X-M_{\boldsymbol{n}}-b,
\]
which is equivalent to \eqref{eq:maximal-component-four}.
\end{proof}

The reduction theorem \cite[Theorem~1.1]{PortillaReduction} gives
$m_{\boldsymbol{n}}\in\{2,3,4\}$.

\section{Common preparation and auxiliary lemmas}
\label{sec:auxiliary}

We choose coordinates, bound weighted orders, and prove a criterion
for noninjectivity, following the three steps of the
\hyperref[sec:strategy]{strategy}.  A smooth component
of the critical locus with no common tangent with $V$ provides the axis
required in Step~\ref{step:ramification}.  Taking this component as an
axis expresses the divisibility of the Jacobian minors as divisibility
by a power of its defining coordinate.  The base points give lower bounds for the orders of the
coefficients of a Weierstrass polynomial defining $V$, whose roots
determine the weights for Step~\ref{step:root-separation}.
\Cref{prop:common-preparation} establishes these coordinates before
we introduce the weighted estimates.  For
Step~\ref{step:third-coordinate}, \Cref{lem:weighted-lifting} proves that
an isolated collision of the initial forms implies noninjectivity
of the analytic map.

\subsection{Base points and branch counts}

We call $p$ the \emph{predecessor} of $q$ if $q$ lies on the
exceptional curve immediately after blowing up $p$.
An infinitely near point $q$ is \emph{proximate} to $p$, written
$q\to p$, if it lies on the strict transform of the exceptional
curve created by blowing up $p$.  A point other than the origin is \emph{free} if it
is proximate to one point, and \emph{satellite} if it is proximate to two.
For a base point $p\in\mathcal C$, put
\[
 \rho_p=m_{\boldsymbol{n},p}
 -\sum_{q\to p}m_{\boldsymbol{n},q}.
\]
The proximity inequalities give $\rho_p\geq0$
\cite[Corollary~7.2.14]{CasasAlvero}.  For a curve germ $\Gamma$,
its \emph{path through the cluster} is the set of points $p\in\mathcal C$
through which its successive strict transforms pass.  For $1\leq j\leq m_{\boldsymbol n}$, put
\begin{equation}
 A_j=\#\{p\in\mathcal C:m_{\boldsymbol n,p}=j\}.
 \label{eq:base-point-counts}
\end{equation}
For coordinates $(x,y)$ and a positive rational number $\alpha$,
define the weighted order
\[
 \nu_\alpha\left(\sum_{i,j\geq0}h_{ij}x^iy^j\right)
 =\min_{h_{ij}\neq0}(\alpha i+j),\qquad \nu_\alpha(0)=\infty.
\]
The next lemma counts the branches of $V$, describes the base points
$p\in\mathcal C$ with $m_{\boldsymbol n,p}=m_{\boldsymbol n}$, and
computes intersection multiplicities with smooth curves.

\begin{lemma}
\label{lem:cluster}
One has
\begin{equation}
 b=m_{\boldsymbol{n}}
 -\sum_{\substack{q\in\mathcal C\\q\text{ satellite}}}
  m_{\boldsymbol{n},q}.                               \label{eq:branch-count}
\end{equation}
The base points satisfying $m_{\boldsymbol n,p}=m_{\boldsymbol n}$
can be ordered as $p_0=0,p_1,\ldots,p_r$, where each $p_{j+1}$ lies
on the exceptional curve created by blowing up $p_j$ and on no other
exceptional curve.  There are coordinates in which
\begin{equation}
\nu_{A_{m_{\boldsymbol n}}}(h)\geq m_{\boldsymbol n}A_{m_{\boldsymbol n}}\qquad(h\in I).         \label{eq:chain-value}
\end{equation}
Let $\Gamma$ be a smooth branch having no common component with
$V$, and suppose that their strict transforms after principalization
of $I$ do not meet over the origin.  Then
\begin{equation}
 i_0(v,\Gamma)=m_{\boldsymbol{n}}+
 \sum_{\substack{p\neq0\text{ on the path}\\
                  \text{followed by }\Gamma}}
                  m_{\boldsymbol{n},p}.                \label{eq:path-intersection}
\end{equation}
Every point on this path other than the origin is free.
\end{lemma}

\begin{proof}
After principalization of $I$, the intersection number of the strict
transform of $V$ with the strict transform of the exceptional curve
created at $p$ is $\rho_p$.  Our choice
of $v$ in \Cref{sec:input} makes these intersections transverse and
disjoint from crossings of the exceptional divisor.  Hence $b=\sum_p\rho_p$.
In this sum, each $m_{\boldsymbol n,q}$ with $q\neq0$ is subtracted
once at the predecessor of $q$, and once more if $q$ is satellite.
This proves \eqref{eq:branch-count}.

The integers $m_{\boldsymbol n,p}$ do not increase along a sequence
of blowups.  The proximity inequality at $p$ permits at most one
point $q$ with predecessor $p$ and $m_{\boldsymbol n,q}=m_{\boldsymbol n}$.
If such a point exists, then $m_{\boldsymbol n,p}=m_{\boldsymbol n}$ as well.
This point $q$ cannot be satellite.  Otherwise, both $q$ and its
predecessor $p$ would be proximate to an earlier point $p'$, whose
proximity inequality would give
$m_{\boldsymbol n,p'}\geq m_{\boldsymbol n,p}+m_{\boldsymbol n,q}=2m_{\boldsymbol n}$.
This contradicts $m_{\boldsymbol n,p'}\leq m_{\boldsymbol n}$.
Thus the base points $p$ with $m_{\boldsymbol n,p}=m_{\boldsymbol n}$
can be listed as $p_0=0,p_1,\ldots,p_r$, with $p_j$ the predecessor
of $p_{j+1}$ and each $p_{j+1}$ free.
Choose coordinates $(x,y)$ so that the successive strict transforms
of the smooth curve $x=0$ pass through $p_0,\ldots,p_r$ and $y=0$ is
transverse to $x=0$.  The successive blowup charts give
$y^{-m_{\boldsymbol n}A_{m_{\boldsymbol n}}}h(y^{A_{m_{\boldsymbol n}}}X,y)\in\C\{X,y\}$ for every
$h\in I$, proving \eqref{eq:chain-value}.  Finally, the blowup
formula for intersection multiplicity, iterated along the smooth
branch $\Gamma$, gives \eqref{eq:path-intersection}.  The intersection
contribution from the strict transforms of $V$ and $\Gamma$ after
principalization vanishes by the disjointness assumption.
After each blowup on the smooth branch $\Gamma$, its strict transform
meets the new exceptional curve transversely away from the older
exceptional curves.  Hence every point on the path of $\Gamma$ other than the origin
is free.
\end{proof}

\subsection{A common preparation for the three orders}

The ramification bound \eqref{eq:one-polar} controls how many points
of a nearby fiber $V_s$ of $v$ can lie on a component $\Gamma_i$
with $m_i=m_{\mathrm{tr}}$.  We now
use the cluster formulas to show that this component is smooth and
meets $V$ with intersection multiplicity $m_{\boldsymbol n}$.
This completes Step~\ref{step:ramification} and provides the
coordinates for the root separation in Step~\ref{step:root-separation}.

We call a monic polynomial of degree $m_{\boldsymbol n}$ in $x$
\emph{centered} if its coefficient of $x^{m_{\boldsymbol n}-1}$
vanishes.  Equivalently, its roots, counted with multiplicity, have
mean zero.  The next proposition produces a centered Weierstrass
polynomial whose zero set is $V$.

\begin{proposition}[Common preparation]
\label{prop:common-preparation}
\label{lem:maximal-chain-weierstrass}
Under the hypotheses of \Cref{sec:input}, choose $i$ with
$m_i=m_{\mathrm{tr}}$.  Then $\Gamma_i$ is smooth,
$a_i=m_{\boldsymbol n}$, and its path through the base-point cluster
consists only of the origin.
There are coordinates $(x,y)$ with $\Gamma_i=\{y=0\}$ in which
\begin{equation}
 \begin{gathered}
 v=EG,\qquad
 G=x^{m_{\boldsymbol n}}+
   \sum_{j=0}^{m_{\boldsymbol n}-2}c_j(y)x^j,\qquad
 E\in\C\{x,y\}^*,\\
 c_j\in\C\{y\},\qquad
 \ord_yc_j\geq(m_{\boldsymbol n}-j)A_{m_{\boldsymbol n}}
 \quad(0\leq j\leq m_{\boldsymbol n}-2).
 \end{gathered}
 \label{eq:maximal-chain-coefficients}
\end{equation}
If $m_{\boldsymbol n}=2$ and $M_{\boldsymbol n}=2$, then
$m_{\mathrm{tr}}=2$.
\end{proposition}

\begin{proof}
Write $r=\mult_0(\Gamma_i)$.  The local intersection inequality,
\eqref{eq:one-polar}, and \eqref{eq:input-knot} give
\[
 m_{\boldsymbol n}r\leq a_i,\qquad
 a_i(m_{\mathrm{tr}}-1)\leq2m_X-M_{\boldsymbol n}-b,\qquad
 M_{\boldsymbol n}\leq m_{\mathrm{tr}}.
\]
If $m_{\boldsymbol n}=4$, \Cref{lem:maximal-component-estimate}
gives $a_i\leq4$, hence $r=1$ and $a_i=4$.

Suppose $m_{\boldsymbol n}=3$.  The same lemma gives
$r=1$ and $a_i\in\{3,4\}$.  If $a_i=4$,
\eqref{eq:path-intersection} requires one further base point $p$ on
the path of $\Gamma_i$, with $p$ free and $m_{\boldsymbol n,p}=1$.  This point is distinct from
the satellite points counted in \eqref{eq:branch-count}, so
\[
 3M_{\boldsymbol n}-m_X=2(A_2+A_1)\geq2,\qquad
 M_{\boldsymbol n}\geq3+1+(3-b)=7-b.
\]
Consequently
\[
 2m_X-7M_{\boldsymbol n}-3b+12
 \leq8-M_{\boldsymbol n}-3b\leq1-2b<0,
\]
contrary to \eqref{eq:maximal-component-four}.  Thus $a_i=3$.

Now suppose $m_{\boldsymbol n}=2$.  The cluster formulas
\eqref{eq:input-point-basis} and \eqref{eq:branch-count} give
\[
 M_{\boldsymbol n}=2A_2+A_1,\qquad
 m_X=2M_{\boldsymbol n}-A_1,\qquad 1\leq b\leq2.
\]
If $M_{\boldsymbol n}\geq3$ and $a_i\geq4$, then
\[
 4(M_{\boldsymbol n}-1)
 \leq a_i(m_{\mathrm{tr}}-1)
 \leq3M_{\boldsymbol n}-b,
\]
forcing $M_{\boldsymbol n}=3$ and $b=1$.  But a cluster with
$M_{\boldsymbol n}=3$ consists of the origin, where
$m_{\boldsymbol n,0}=2$, and one point $p$ on the exceptional curve
of the first blowup, with
$m_{\boldsymbol n,p}=1$.  Thus \eqref{eq:branch-count} gives $b=2$.
Therefore $a_i\leq3$ and $r=1$.
If $a_i=3$, the path of $\Gamma_i$ contains a free base point $p$
with $m_{\boldsymbol n,p}=1$.  By \eqref{eq:branch-count}, the satellite
points account for $2-b$ of the $A_1$ points.  Including this free
point gives $A_1\geq3-b$.  The polar bound \eqref{eq:one-polar} would then imply
\[
 3(M_{\boldsymbol n}-1)
 \leq3M_{\boldsymbol n}-2A_1-b,\qquad
 3\geq2A_1+b\geq6-b\geq4.
\]
Hence $a_i=2$ when $M_{\boldsymbol n}\geq3$.

If $M_{\boldsymbol n}=2$, the origin is the only base point, so
$m_X=4$ and $b=2$.  The quadratic initial form $q_v$ of $v$ has two
distinct roots in $\mathbb P^1$.  Since $i_0(v,u)=4$, the initial form
$q_u$ of $u$ is also quadratic and has no common root with $q_v$.
The morphism
\[
 [q_v:q_u]:\mathbb P^1\longrightarrow\mathbb P^1,\qquad
 [x:y]\longmapsto[q_v(x,y):q_u(x,y)]
\]
has degree two.  Riemann--Hurwitz gives two simple ramification
points, neither at a zero of $q_v$, since $q_u/q_v$ has simple
poles there.  Thus $J(q_v,q_u)$ is squarefree and coprime to $q_v$.
Let $h_i=0$ be an irreducible equation for $\Gamma_i$.  The initial
form of $h_i$ for total degree is a nonzero multiple
of $L^r$ for a linear form $L$.  The divisibility
$h_i^{m_{\mathrm{tr}}-1}\mid J(v,u)$ from \eqref{eq:input-divisor} gives
\[
 L^{r(m_{\mathrm{tr}}-1)}\mid J(q_v,q_u).
\]
It follows that $r=1$ and $m_{\mathrm{tr}}=2$.  Since $L$ does not
divide $q_v$, the curves $\Gamma_i$ and $V$ have no common tangent,
and $a_i=2$.

Thus $\Gamma_i$ is smooth and $a_i=m_{\boldsymbol n}$ in every
order.  Formula \eqref{eq:path-intersection} then excludes later
base points on the path of $\Gamma_i$.
Choose coordinates $(x,y)$ with $\Gamma_i=\{y=0\}$ so that the
successive strict transforms of $x=0$ pass through the base points
with $m_{\boldsymbol n,p}=m_{\boldsymbol n}$, ordered as in
\Cref{lem:cluster}.  These axes are transverse.
By \eqref{eq:chain-value}, Weierstrass preparation
\cite[Theorem~I.1.6]{GreuelLossenShustin} writes $v$ as a unit times
\[
 x^{m_{\boldsymbol n}}+\sum_{j<m_{\boldsymbol n}}d_j(y)x^j,
 \qquad
 \ord_y d_j\geq(m_{\boldsymbol n}-j)A_{m_{\boldsymbol n}}.
\]
The translation $x'=x+d_{m_{\boldsymbol n}-1}(y)/m_{\boldsymbol n}$
centers the roots of this Weierstrass polynomial.  The added term
$d_{m_{\boldsymbol n}-1}(y)/m_{\boldsymbol n}$ has order at least
$A_{m_{\boldsymbol n}}$, so it preserves the coefficient bounds and
the axis $\Gamma_i$.  This proves \eqref{eq:maximal-chain-coefficients}.
\end{proof}

\subsection{The first separation of the roots}

For the formulas in this subsection, write $m=m_{\boldsymbol n}$.
The centered polynomials supplied by \Cref{prop:common-preparation}
have the following forms in the three orders under consideration:
\[
 \begin{array}{c|c}
 m&G(x,y)\\ \hline
 2&x^2+c_0(y)\\
 3&x^3+c_1(y)x+c_0(y)\\
 4&x^4+c_2(y)x^2+c_1(y)x+c_0(y)
 \end{array}.
\]
For Step~\ref{step:root-separation}, we divide the roots by powers of
$y$ and find the smallest exponent for which their limits differ as
$y\to0$.  We use weighted initial forms to compute these limits and
to compare the three coordinates $v,u,w$ of $\boldsymbol n$.
For positive integral weights $(w_x,w_y)$, the \emph{weighted
order} of a nonzero series is the least value of $iw_x+jw_y$
among its monomials $x^iy^j$.  Its \emph{weighted initial form},
or simply \emph{initial form}, is the sum of the monomials of that
weight.  We use
$\operatorname{ord}_{(w_x,w_y)}$ and
$\operatorname{in}_{(w_x,w_y)}$ for these operations.  The same
definitions apply to positive rational weights, which may be multiplied
by a common denominator.
A polynomial is \emph{weighted homogeneous} if all its nonzero monomials
have the same weight.

For a reduced Weierstrass polynomial $G$ as in
\eqref{eq:maximal-chain-coefficients}, define its \emph{separation exponent} by
\begin{equation}
 \alpha=\min_{0\leq j\leq m-2}\frac{\ord_yc_j}{m-j}\geq A_m,
 \qquad \ord_y0=\infty.
 \label{eq:separation-weight}
\end{equation}
At least one $c_j$ is nonzero, since otherwise $G=x^m$ would not
be reduced for $m\geq2$.  Thus $\alpha$ is finite.  By Puiseux' theorem
\cite[Theorem~2.2.6]{Wall}, after a substitution $y=t^d$ for some
integer $d\geq1$, the roots of $G(x,y)=0$ are
convergent power series in $t$.  Denote the corresponding series in
fractional powers of $y$ by $\xi_1(y),\ldots,\xi_m(y)$.
These are the \emph{Puiseux roots} of $G$.  Normalize their orders by
$\ord_y y=1$.  Each $\xi_i$ has order at least $\alpha$:
otherwise its $m$th power would be the unique term of least order
in $G(\xi_i(y),y)=0$.  As $y\to0$, the substitution $x=y^\alpha z$ gives the
\emph{limiting polynomial} $P_m$, defined by
\begin{equation}
 y^{-m\alpha}G(y^\alpha z,y)\longrightarrow P_m(z),\qquad
 \operatorname{in}_{(\alpha,1)}G=y^{m\alpha}P_m(x/y^\alpha).
 \label{eq:limiting-polynomial}
\end{equation}
The limits of $\xi_i(y)/y^\alpha$ are the roots of the monic polynomial
$P_m$, counted with multiplicity.  Its coefficient of $z^{m-1}$
vanishes, and the minimum in \eqref{eq:separation-weight} gives a
nonzero coefficient of lower degree.  If $P_m(z)=(z-\zeta)^m$ for some
$\zeta\in\C$, its vanishing coefficient of $z^{m-1}$ would give
$\zeta=0$, contrary to the existence of that nonzero coefficient
of lower degree.
Since $P_m$ is monic of degree $m$, it is therefore not the $m$th
power of a polynomial in $\C[z]$, and its roots are not all equal.
For every rational number $\beta$ with $0<\beta<\alpha$, all
rescaled roots $\xi_i(y)/y^\beta$ tend to zero as $y\to0$.  Thus
$\alpha$ is the first exponent at which at least two rescaled roots
have different limits.

\subsection{Collisions of weighted initial forms}

To detect noninjectivity, we study pairs of distinct points with the
same image under the map formed by the weighted initial forms of
$v,u,w$.  These are the pairs whose
separation is required in Step~\ref{step:third-coordinate}.
Weighted homogeneity allows us to record them in a weighted
projective space.  The lemma below shows that an isolated collision
in that space implies that the analytic map is not injective.

Use positive integral weights $(w_x,w_y)$, clearing denominators
if the given weights are rational.  The weighted projective space
$\mathbb P(w_x,w_y,w_x,w_y)$ identifies nonzero tuples $(x,y,X,Y)$
under the rescalings
\[
 (x,y,X,Y)\longmapsto
 (\lambda^{w_x}x,\lambda^{w_y}y,\lambda^{w_x}X,\lambda^{w_y}Y),
 \qquad \lambda\in\C^*.
\]
If
$F_0=(F_{1,0},F_{2,0},F_{3,0}):\C^2\to\C^3$ has nonzero weighted
homogeneous polynomial components for these weights, let
$\mathcal Z(F_0)$ be the closed algebraic set in the weighted projective space $\mathbb P(w_x,w_y,w_x,w_y)$ cut out by
\[
 F_{i,0}(x,y)-F_{i,0}(X,Y)=0
 \qquad(1\leq i\leq3).
\]
Its \emph{diagonal} is
\[
 \Delta=V(x-X,y-Y)\subset\mathbb P(w_x,w_y,w_x,w_y).
\]
When only two weighted homogeneous forms $P,Q$ are specified,
$\mathcal Z(P,Q)$ denotes the closed algebraic set defined by
$P(x,y)=P(X,Y)$ and $Q(x,y)=Q(X,Y)$ in the same weighted projective
space.  We call this set the \emph{projective fiber product} of
$(P,Q)$ with itself.  For a third form $R$, write
\[
 \Delta_R=V\bigl(R(x,y)-R(X,Y)\bigr).
\]
Thus $\mathcal Z(P,Q,R)=\mathcal Z(P,Q)\cap\Delta_R$.
We call a point of $\mathcal Z(P,Q,R)$ outside $\Delta$ a
\emph{collision of the initial forms}.  We will use the following criterion to pass from such
a point to noninjectivity of an analytic germ.
For an analytic germ $F=(F_1,F_2,F_3):(\C^2,0)\to(\C^3,0)$
with $F_i\neq0$, let $F_0$ be the map whose components are the
initial forms of the $F_i$ for the fixed weights $(w_x,w_y)$.

\begin{lemma}
\label{lem:weighted-lifting}
If $\mathcal Z(F_0)$ has an isolated point outside $\Delta$, then $F$
is not injective.
\end{lemma}

\begin{proof}
Choose a homogeneous coordinate of weight $d$ that is nonzero at the isolated point
of $\mathcal Z(F_0)\setminus\Delta$.  Setting this coordinate equal to
one gives a finite map from $\C^3$ onto the corresponding chart of
$\mathbb P(w_x,w_y,w_x,w_y)$. Indeed, this map is the quotient by the
residual action of the finite group of $d$th roots of unity.
Choose a preimage of the isolated point
and local coordinates $\xi_1,\xi_2,\xi_3$ at that preimage.
Under weighted rescaling, the equations
\[
 \tau^{-\operatorname{ord}_{(w_x,w_y)}F_i}
 \bigl(F_i(\tau^{w_x}x,\tau^{w_y}y)
       -F_i(\tau^{w_x}X,\tau^{w_y}Y)\bigr)=0
 \qquad(1\leq i\leq3)
\]
have left-hand sides that are convergent in
$(\xi_1,\xi_2,\xi_3,\tau)$ and have an isolated common zero on
$\{\tau=0\}$.  Three equations in four
variables have a common zero set of dimension at least one
\cite[Theorem~B.2.1(1)]{GreuelLossenShustin}.  A component of positive dimension through the chosen zero is
therefore not contained in $\{\tau=0\}$.
For $\tau\neq0$, rescaling gives two distinct points of $\C^2$
with the same image under $F$.  Both points approach the origin
as $\tau\to0$, proving noninjectivity.
\end{proof}

\subsection{Weighted orders and divisibility of the Jacobian minors}
\label{sec:weighted-pair}

We compare the separation exponent of $G$ with the vanishing order
of $J(v,u)$ along $\Gamma_i$.  The following estimate uses
\eqref{eq:input-divisor} and \eqref{eq:input-knot} to connect
Steps~\ref{step:ramification} and \ref{step:root-separation}.
It applies when $m_{\boldsymbol n}\in\{3,4\}$.
Retain the coordinates of \Cref{prop:common-preparation}, put
$m=m_{\boldsymbol n}\in\{3,4\}$, and use $\alpha,P=P_m$ from
\eqref{eq:separation-weight}--\eqref{eq:limiting-polynomial}.
By \eqref{eq:input-divisor}, $y^{m_{\mathrm{tr}}-1}\mid J(v,u)$.
Throughout this subsection, initial forms and weighted orders are
taken with respect to $(\alpha,1)$.
Subtract terms $cv^q$ from $u$ while its initial form is a scalar
multiple of a positive integral power of the initial form of $v$.
These target changes leave $u|_V$ unchanged and preserve the ideals
$I$ and $(v,u)B$ and the Jacobian determinant $J(v,u)$.
Each subtraction strictly increases the weighted order of the remainder.
If the process did not terminate, its remainders would have
arbitrarily large order at the origin, so $u\in(v)+\mathfrak m^N$
for every $N$.  Krull intersection in $B/(v)$
\cite[Theorem~B.4.2]{GreuelLossenShustin} would give $u\in(v)$,
contradicting $i_0(v,u)=m_X<\infty$ in \eqref{eq:input-point}.
For this choice of $u$, put
\[
 \ell=\nu_\alpha(u),\qquad
 \operatorname{in}_{(\alpha,1)}u=y^\ell R(x/y^\alpha),\qquad
 k=\deg R.
\]
Define the polynomial $\Omega=\ell P'R-m\alpha PR'$.
If $\Omega\neq0$, set $j=\deg\Omega$.

\begin{lemma}
\label{lem:weighted-pair}
One has $\ell=k\alpha$ and
\begin{equation}
 m_X\geq m_{\boldsymbol{n}}\ell.                      \label{eq:resultant-bound}
\end{equation}
If $\Omega\neq0$, then
\begin{equation}
 m_{\mathrm{tr}}
 \leq\ell+\alpha(m_{\boldsymbol{n}}-1-j).              \label{eq:jacobian-bound}
\end{equation}
In this case, $j\leq m_{\boldsymbol n}-2$ and
$k\geq m_{\boldsymbol n}+1$.  If $\Omega=0$, then
\begin{equation}
 R^{m_{\boldsymbol{n}}}=\kappa P^k                    \label{eq:power-identity}
\end{equation}
for some $\kappa\in\C^*$, and the inequality in
\eqref{eq:resultant-bound} is strict.
\end{lemma}

\begin{proof}
Every Puiseux root $\xi(y)$ of $G$ has order at least $\alpha$,
as established before \eqref{eq:limiting-polynomial}.
By \eqref{eq:input-point}, the intersection formula
\cite[Lemma~1.2.1(ii)--(iii)]{Wall}, with roots counted with multiplicity and
$\ord_y y=1$, gives
\[
 m_X=i_0(G,u)=\sum_{G(\xi,y)=0}\ord_y u(\xi(y),y)\geq m\ell.
\]
If $\Omega\neq0$, differentiation gives
\[
 \operatorname{in}_{(\alpha,1)}J(v,u)
 =E(0,0)y^{\ell+(m-1)\alpha-1}\Omega(x/y^\alpha).
\]
Indeed, in $J(v,u)=EJ(G,u)+GJ(E,u)$, the summand $GJ(E,u)$ has
weighted order at least
$\ell+(m-1)\alpha-1+\min(\alpha,1)$, so it does not contribute
to the initial form of $J(v,u)$.
The nonzero coefficient of $z^j$ in $\Omega(z)$ gives a nonzero
monomial $x^jy^{\ell+(m-1-j)\alpha-1}$ in that initial form.
Since $y^{m_{\mathrm{tr}}-1}\mid J(v,u)$, its exponent of $y$
is at least $m_{\mathrm{tr}}-1$, proving \eqref{eq:jacobian-bound}.
Combining \eqref{eq:jacobian-bound}
with \eqref{eq:resultant-bound} and \eqref{eq:input-knot} yields
\[
 (m-3)\ell+M_{\boldsymbol n}\leq3\alpha(m-1-j).
\]
Since $m\geq3$, $\ell>0$, and $M_{\boldsymbol n}\geq m>0$,
the quantity $(m-3)\ell+M_{\boldsymbol n}$ is positive.
As $\alpha>0$, this forces $m-1-j>0$, hence $j\leq m-2$.
The coefficient of $z^{m+k-1}$ in $\Omega(z)$ is
$m(\ell-k\alpha)$ times the nonzero coefficient of $z^k$ in $R(z)$.
Since $m+k-1\geq m-1>j=\deg\Omega$, this coefficient vanishes,
so $\ell=k\alpha$.
Thus $x^k$ occurs in $u$, and $u\in\mathfrak m^m$ gives $k\geq m$.
If $k=m$, let $r_m\neq0$ be the leading coefficient of $R$.
Since $P$ is monic and the choice of $u$ in \Cref{sec:weighted-pair} excludes
$R=r_mP$, the polynomial $\widetilde R=R-r_mP$ is nonzero and
has degree $d<m$.  The equality $\ell=m\alpha$ gives
\[
 \Omega=m\alpha(P'\widetilde R-P\widetilde R').
\]
If $\widetilde r_d$ is the leading coefficient of $\widetilde R$,
the coefficient of $z^{m+d-1}$ in this expression is
$m\alpha(m-d)\widetilde r_d\neq0$.
Hence $\deg\Omega=m+d-1\geq m-1$, contradicting $j\leq m-2$.

If $\Omega=0$, the coefficient of $z^{m+k-1}$ in the zero
polynomial $\ell P'R-m\alpha PR'$ gives $\ell=k\alpha>0$.
The identity $kP'R-mPR'=0$ then implies
$(R^m/P^k)'=0$ in the rational function field $\C(z)$.
A rational function with zero derivative is constant, so
$R^m/P^k=\kappa\in\C^*$, proving \eqref{eq:power-identity}.
For every Puiseux root $\xi$ of $G$, the constant term
$\zeta$ of $\xi/y^\alpha$ satisfies $P(\zeta)=0$, and
therefore $R(\zeta)=0$.  Each order $\ord_y u(\xi(y),y)$ in the
resultant formula is then strictly greater than $\ell$.
\end{proof}

\section{Exclusion of
\texorpdfstring{$m_{\boldsymbol{n}}=4$}{m(n)=4}}
\label{sec:dfour}

When $m_{\boldsymbol n}=4$, we compare the separation of the roots of the
Weierstrass polynomial $G$ defining $V$ with the multiplicities of the
successive strict transforms of $V$.
In Step~\ref{step:root-separation}, the first separation must be into
two pairs, and the multiplicities of the strict transforms of $V$
after the corresponding blowups violate the numerical bounds.  This excludes $m_{\boldsymbol n}=4$ without the
argument for the third coordinate in Step~\ref{step:third-coordinate}.

\begin{theorem}
\label{thm:dfour}
Under the hypotheses of \Cref{sec:input}, one has
$m_{\boldsymbol{n}}\neq4$.
\end{theorem}

\begin{proof}
Assume $m_{\boldsymbol n}=4$ and apply \Cref{prop:common-preparation},
with $v=EG$.  Let $\alpha\geq A_4$ be the separation exponent
from \eqref{eq:separation-weight}, and write $P=P_4$ for the
limiting quartic in \eqref{eq:limiting-polynomial}.
Since $a_i=4$, \eqref{eq:maximal-component-four} holds.

To determine $P$, subtract scalar multiples of powers of $v$ from $u$
as in \Cref{sec:weighted-pair}, until the initial form of $u$ for
weights $(\alpha,1)$ is not a scalar multiple of a positive integral
power of the initial form of $v$.  Use the notation $R,\Omega$ and
$k=\deg R$ from \Cref{lem:weighted-pair}.
If $\Omega\neq0$, the lemma gives
$k\geq5$, $m_X\geq4k\alpha$, and $m_{\mathrm{tr}}\leq(k+3)\alpha$.
Also, \eqref{eq:input-resolution} gives $M_{\boldsymbol n}\geq m_X/4$, so
\[
 5k\alpha\leq\frac54m_X\leq m_X+M_{\boldsymbol n}
 \leq3m_{\mathrm{tr}}\leq3(k+3)\alpha,
\]
contradicting $k\geq5$.

Hence $\Omega=0$, and \eqref{eq:power-identity} gives
$R(z)^4=\kappa P(z)^k$ in $\C[z]$, with $\kappa\neq0$.
Factoring both sides into linear factors shows that $P$ and $R$
have the same roots and, at each such root $\zeta$,
\[
 4\ord_{z=\zeta}R=k\ord_{z=\zeta}P,
\]
where $\ord_{z=\zeta}$ denotes the order of vanishing at $\zeta$.
If $k$ is odd, every root multiplicity of $P$ is divisible by
four, contrary to the fact that $P$ is not a fourth power, established
after \eqref{eq:limiting-polynomial}.
If $4\mid k$, the identity $R^4=\kappa P^k$ gives $R=cP^{k/4}$, contrary to the
choice of $u$ in \Cref{sec:weighted-pair}.  Thus $k\equiv2\pmod4$, and every root
multiplicity of $P$ is even.  Write $P=Q^2$ with $Q$ a monic quadratic.
The vanishing cubic coefficient of $P$ gives
\[
 Q=z^2+\eta,\qquad P=(z^2+\eta)^2,\qquad \eta\neq0,
\]
where $\eta\neq0$ because $P$ is not a fourth power.
Since $x^k$ occurs in $u\in\mathfrak m^4$, we have $k\geq4$,
and hence $k\geq6$.
The roots of $G$, divided by $y^\alpha$, tend in pairs to $a,-a$
with $a^2=-\eta$.
The strict inequality in \eqref{eq:resultant-bound} when $\Omega=0$ gives
\begin{equation}
 m_X>4k\alpha\geq24A_4.                         \label{eq:d4-contact}
\end{equation}


We now determine the multiplicities of the strict transforms of $V$
after this separation.
The coefficient $c_2\in\C\{y\}$ of $x^2$ in $G$ has leading term
$2\eta y^{2\alpha}$.  Since $\eta\neq0$, $2\alpha=\ord_yc_2$ is an integer.
Since $\alpha\geq A_4$, the number $q=2(\alpha-A_4)$ is a
nonnegative integer.  By \Cref{lem:cluster}, the $A_4$ base points $p$
with $m_{\boldsymbol n,p}=4$ form a sequence in which each point after
the origin lies on the exceptional curve created by blowing up the
preceding point.  After blowing up these points in order, the strict
transform of $V$ is defined by the polynomial
\[
 G_{A_4}(X,y)=y^{-4A_4}G(y^{A_4}X,y)
 =X^4+\widetilde c_2(y)X^2+\widetilde c_1(y)X+\widetilde c_0(y),
\]
whose coefficients satisfy
\[
 \ord_y\widetilde c_2=q,\qquad
 \ord_y\widetilde c_1>\tfrac32q,\qquad
 \ord_y\widetilde c_0=2q.
\]
If $q\geq2$, every monomial of $G_{A_4}$ has ordinary total
degree at least four, and $X^4$ occurs with coefficient one.
Thus the strict transform would have multiplicity exactly four at
$(0,0)$.  By \Cref{sec:input}, this would be a further base point $p$
with $m_{\boldsymbol n,p}=4$, contradicting the definition of $A_4$
in \eqref{eq:base-point-counts}.
If $q=0$, the strict transform meets the exceptional curve $y=0$
at $(a,0)$ and $(-a,0)$, since
$G_{A_4}(X,0)=(X^2+\eta)^2$ and $a^2=-\eta\neq0$.
At either point, the restriction of $G_{A_4}$ to $y=0$ has order two in
$X-a$ or $X+a$, respectively.  The multiplicity of the plane
curve germ $\{G_{A_4}=0\}$ at that point is therefore at
most two.
If $q=1$, the restriction $G_{A_4}(X,0)=X^4$ gives only the
point $(0,0)$ on the exceptional curve.  The lowest homogeneous
part of $G_{A_4}$ for the ordinary total degree is
$\eta^2y^2\neq0$, so the plane curve germ has multiplicity two there.
Every remaining base point maps to one of the points just described
under the intervening blowups.
Every Puiseux root $\xi(y)$ of $G$ has order at least $\alpha\geq A_4$.
Hence $\xi(y)/y^r$ remains bounded for $1\leq r\leq A_4$ and
tends to zero for $r<A_4$.  Thus the successive strict transforms of
every branch of $V$ lie in the charts with coordinates $(x/y^r,y)$,
and after blowing up all $A_4$ points they pass through one of the
points just listed.
By \cite[Lemma~3.4.2]{Wall} and the local intersection inequality,
multiplicity cannot increase under further blowups.  Since the
multiplicities of the strict transforms of $V$ at the base points are
$m_{\boldsymbol n,p}$, all remaining base points satisfy
$m_{\boldsymbol n,p}\leq2$.  Therefore
\eqref{eq:input-point-basis} gives
\begin{equation}
 m_X\leq16A_4+2(M_{\boldsymbol n}-4A_4)
      =2M_{\boldsymbol n}+8A_4.                    \label{eq:d4-two}
\end{equation}
Combining \eqref{eq:d4-two} with \eqref{eq:d4-contact} gives
$M_{\boldsymbol n}>8A_4\geq8$ and $m_X<3M_{\boldsymbol n}$.
Thus the average $m_X/M_{\boldsymbol n}$ of the multiplicities $m_{\boldsymbol n,p}$,
with weights proportional to $m_{\boldsymbol n,p}$, is below three.
However, integrality gives $M_{\boldsymbol n}\geq9$, and the polar
estimate \eqref{eq:maximal-component-four} gives
\[
 2m_X\geq7M_{\boldsymbol n}+3b-12
 \geq7M_{\boldsymbol n}-9\geq6M_{\boldsymbol n},
\]
a contradiction.
\end{proof}

\section{Exclusion of
\texorpdfstring{$m_{\boldsymbol{n}}=3$}{m(n)=3}}
\label{sec:dthree}

When $m_{\boldsymbol n}=3$, the numerical bounds determine two possible
pairs of weighted initial forms of $v,u$ in Step~\ref{step:root-separation}.
For Step~\ref{step:third-coordinate}, the projective fiber product
records pairs of points with the same image under the map defined by
each pair of forms, up to simultaneous rescaling.
The divisibility conditions on the Jacobian minors prevent the third
coordinate from separating all such pairs.  The isolated collisions
obtained from these initial forms lift to the analytic map $\boldsymbol n$
and contradict
\eqref{eq:strategy-separation}.

In the next two subsections, assume $m_{\boldsymbol n}=3$ and use
the component $\Gamma_i$ and coordinates of \Cref{prop:common-preparation}.
Thus $\Gamma_i=\{y=0\}$ and $a_i=3$.  Put $a=A_3\geq1$.
Since $m_{\boldsymbol n,p}\in\{1,2,3\}$ at every base point $p\in\mathcal C$,
\eqref{eq:input-point-basis} and \eqref{eq:branch-count} give
\begin{equation}
 2M_{\boldsymbol n}-m_X
 =\sum_{p\in\mathcal C}m_{\boldsymbol n,p}(2-m_{\boldsymbol n,p})
 \geq-3a,\qquad b\leq3.
 \label{eq:d3-point-bounds}
\end{equation}
\subsection{The two possible pairs of weighted initial forms}

In independent variables $x,t$, define the pairs of homogeneous polynomials
\begin{align}
 E:\quad &P=x^3+3t^3,& Q&=x^4+4xt^3,               \label{eq:d3-E}\\
 H:\quad &P=x^3+3xt^2,& Q&=x^4+4x^2t^2+2t^4.       \label{eq:d3-H}
\end{align}
The Jacobian of each pair is
\begin{equation}
 J_{x,t}(P,Q)=
 \begin{cases}-36t^5,&E,\\24t^5,&H.\end{cases}    \label{eq:d3-PQjac}
\end{equation}
The following proposition identifies $E$ and $H$, after substituting
$t=y^a$, as the possible pairs of initial forms of $v,u$ for weights
$(a,1)$ and determines the transverse multiplicity $m_{\mathrm{tr}}$.

\begin{proposition}
\label{prop:d3-models}
Under the assumptions of \Cref{sec:dthree}, there are source
coordinates $(x,y)$ with $\Gamma_i=\{y=0\}$ and
changes of target coordinates, obtained by subtracting a polynomial in
$v$ from $u$ and multiplying $v,u$ by nonzero constants, such that
\[
 \operatorname{in}_{(a,1)}v=P(x,y^a),\qquad
 \operatorname{in}_{(a,1)}u=Q(x,y^a)
\]
for one of the pairs \eqref{eq:d3-E}, \eqref{eq:d3-H}.  Moreover,
\begin{equation}
 m_X=12a,\qquad m_{\mathrm{tr}}=6a.                \label{eq:d3-exact}
\end{equation}
\end{proposition}

\begin{proof}
Let $\alpha\geq a$ be the separation exponent of $G$ from
\eqref{eq:separation-weight}, and write the limiting cubic as
\[
 p=P_3=z^3+a_pz+b_p.
\]
As established after \eqref{eq:limiting-polynomial}, the monic cubic
$p\in\C[z]$ is not equal to $(z-c)^3$ for any
$c\in\C$.  Its root multiplicities sum to three, so at least one
root is simple.
Apply \Cref{lem:weighted-pair} to $(v,u)$ and retain its notation
$R,\Omega,\ell,k$.
If $\Omega=0$, the identity $R^3=\kappa p^k$ in $\C[z]$
gives $3\ord_{z=c}R=k$ at a simple root $c$ of $p$, so $3\mid k$.
At every root of $p$, the multiplicity as a root of $R$ is
$k/3$ times the multiplicity as a root of $p$, and $R$ has no other roots.
Thus $R=c_0p^{k/3}$ for some $c_0\in\C^*$.
Since $\ell=k\alpha$, the initial form of $u$ for weights $(\alpha,1)$ would then be a
scalar multiple of the $(k/3)$rd power of that of $v$, contrary
to the choice of $u$ in \Cref{sec:weighted-pair}.
Thus $\Omega\neq0$.  With $j=\deg\Omega$ and $k=\deg R$, the lemma gives
\[
 j\leq1,\quad k\geq4,\quad \ell=k\alpha,\quad
 m_X\geq3k\alpha,\quad
 m_{\mathrm{tr}}\leq\alpha(k+2-j).
\]
By \eqref{eq:d3-point-bounds} and \eqref{eq:input-knot},
\[
 \frac92k\alpha-\frac32a
 \leq m_X+M_{\boldsymbol n}\leq3m_{\mathrm{tr}}
 \leq3\alpha(k+2-j).
\]
Hence $(k-4+2j)\alpha\leq a\leq\alpha$.
Dividing by $\alpha>0$ gives $k-4+2j\leq1$.
Since $k\geq4$ and $j\geq0$ are integers, this forces
$j=0$.  At every root $c$ of $p$,
$0\neq\Omega(c)=\ell p'(c)R(c)$.
Hence $p'(c)\neq0$ and $R(c)\neq0$ at every root of $p$.
Thus all roots of $p$ are simple, and $p,R$ are coprime in
$\C[z]$.  Along each of the three Puiseux roots of $G$,
the orders of the restrictions of $u$ and $G_x$, normalized by
$\ord_y y=1$, are therefore exactly $k\alpha$ and $2\alpha$,
respectively.  Summing these orders in the
resultant formulas for $i_0(G,u)=m_X$ and $i_0(G,G_x)$ gives
\[
 m_X=3k\alpha,\qquad i_0(G,G_x)=6\alpha.
\]
Since $v=EG$ and $E$ is a unit, $v_x\equiv EG_x\pmod v$
gives $i_0(v,v_x)=i_0(G,G_x)=6\alpha$.
The axis $y=0$ is transverse to $V$, so $i_0(v,y)=3$.
The polar identity
\cite[Lemma~7.2.1 and its consequence, pp.~160--161]{Wall} is
\[
 i_0(v,v_x)=\mu(V,0)+i_0(v,y)-1=\mu(V,0)+2.
\]
Thus $\mu(V,0)=6\alpha-2$, and \eqref{eq:input-milnor} yields
\[
 M_{\boldsymbol n}=3\alpha(k-2)-b+3.
\]
Substituting $m_X=3k\alpha$ and
$M_{\boldsymbol n}=3\alpha(k-2)-b+3$ into
$m_X+M_{\boldsymbol n}\leq3m_{\mathrm{tr}}\leq3\alpha(k+2)$
gives $3\alpha(k-4)\leq b-3$.
Since $\alpha\geq a\geq1$, $k\geq4$, and
\eqref{eq:branch-count} gives $b\leq3$, we have
$0\leq3\alpha(k-4)\leq b-3\leq0$, hence $k=4$ and $b=3$.
Substituting these values into the formulas for $m_X,M_{\boldsymbol n}$
and using $M_{\boldsymbol n}\leq m_{\mathrm{tr}}\leq\alpha(k+2)$ gives
\[
 k=4,\qquad b=3,\qquad m_X=12\alpha,
 \qquad M_{\boldsymbol n}=m_{\mathrm{tr}}=6\alpha.
\]
The multiplicities of the three branches of $V$ sum to $\mult_0V=3$,
so each branch is smooth.  Transversality to $y=0$ gives three convergent
graphs $x=\xi_s(y)$, and the limits of $\xi_s(y)/y^\alpha$ as $y\to0$ are the distinct roots
$c_s$ of $p$.  At least one $c_s$ is nonzero, so the corresponding
series $\xi_s(y)$ has order $\alpha$, proving $\alpha\in\Z$.
For $0\leq r<\alpha$, the strict transforms of all three branches under the
substitutions $x=y^r x_r$ pass through $(x_r,y)=(0,0)$.
After the $\alpha$th blowup, they meet the exceptional curve at
the distinct points $(c_s,0)$.
Thus the strict transforms of all three branches pass through exactly
$\alpha$ common centres of blowup, including the origin.  At each centre
the strict transform of $V$ has multiplicity three, so each centre is a
base point $p$ with $m_{\boldsymbol n,p}=3$.
Therefore $a\geq\alpha$, so $a=\alpha$, proving
\eqref{eq:d3-exact}.

It remains to classify the initial forms of $v,u$ for weights $(a,1)$.
Set $\chi=R$
and write $P=t^3p(x/t)$, $Q=t^4\chi(x/t)$.
Since $\ell=4\alpha$, the nonzero constant $\Omega$ satisfies
$\Omega=\alpha(4p'\chi-3p\chi')$.
Put $c=\Omega/\alpha\neq0$, so $4p'\chi-3p\chi'=c$.
Comparing coefficients of powers of $z$ in this polynomial identity,
with $c_4$ the coefficient of $z^4$ in $\chi(z)$, gives
\[
 \chi=c_4\left(z^4+\frac{4a_p}{3}z^2+\frac{4b_p}{3}z
                      +\frac{2a_p^2}{9}\right),\qquad
 a_pb_p=0,\qquad
 c=c_4\left(\frac{8a_p^3}{9}-4b_p^2\right).
\]
The identity $a_pb_p=0$ implies that at least one of $a_p,b_p$ is zero,
and $c\neq0$ excludes $a_p=b_p=0$.  Thus exactly one is zero.
Rescaling $x,v,u$ makes the nonzero member of $\{a_p,b_p\}$
equal to three and both polynomials $P,Q$ monic in $x$, giving
\eqref{eq:d3-E} or \eqref{eq:d3-H}.
\end{proof}

In these coordinates, \eqref{eq:input-divisor} and
\eqref{eq:d3-exact} imply
\begin{equation}
 y^{6a-1}\mid J(n_j,n_\ell)
 \qquad(1\leq j<\ell\leq3).                          \label{eq:d3-minors}
\end{equation}
The chain rule preserves this divisibility when a polynomial in $v,u$
is subtracted from the third coordinate, as in the proof of
\Cref{thm:dthree}.

\subsection{Collisions of the pairs of homogeneous forms}

For \eqref{eq:d3-E} put $q=3$, and for \eqref{eq:d3-H} put $q=2$.
The transformations $(x,t)\mapsto(x,\xi t)$, with $\xi^q=1$,
define the \emph{cyclic action} used below.
A polynomial $R$ is \emph{invariant} under this action if
$R(x,\xi t)=R(x,t)$ for every $\xi^q=1$.
In coordinates $[x:t:X:T]$ on $\mathbb P^3$, define the lines
\[
 L_\xi=\{(X,T)=(x,\xi t)\}\qquad(\xi^q=1).
\]
The line $L_1$ is the diagonal.
For each pair $(P,Q)$, the projective fiber product $\mathcal Z(P,Q)$
consists of these lines and the inverse image of a curve under the
substitutions $w=t^q$ and $\widetilde w=T^q$.
The points of this fiber product outside the diagonal represent
pairs of distinct points with equal values of $P,Q$, as in
Step~\ref{step:third-coordinate}.
A collision of the three forms $P,Q,R$ may lie on one of the lines $L_\xi$.
If the third form $R$ is invariant under this action, all three forms
are polynomials in $x,w$, where $w=t^q$.
In case $H$, they take equal values at $(0,w),(0,-w)$ with $w\neq0$
when the degree of $R$ in $x,t$ is odd, and at
$(x,-x^2/3),(-x,-x^2/3)$ with $x\neq0$ when that degree is even.
In case $E$, the derivative
identity \eqref{eq:d3-derivative} forces a collision.
For the remaining forms $R$, we compare $R(x,t)-R(X,T)$
with $t^{\deg R}-T^{\deg R}$.
The map defined by $P,Q,t^{\deg R}$ has a collision, and the
comparison prevents all zeros of $R(x,t)-R(X,T)$ on the remaining
curve components of $\mathcal Z(P,Q)$ from lying on the diagonal.  Put
\[
 \Delta=V(x-X,t-T),\qquad p_0=[1:0:1:0],\qquad
 \Delta_R=V\bigl(R(x,t)-R(X,T)\bigr).
\]

\begin{proposition}
\label{prop:d3-cubic-collision}
Let $(P,Q)$ be either pair \eqref{eq:d3-E}, \eqref{eq:d3-H}, and let
$R\in\C[x,t]$ be homogeneous.  If
\begin{equation}
 R\notin\C[P,Q],\qquad
 t^5\mid J_{x,t}(P,R),\qquad t^5\mid J_{x,t}(Q,R),  \label{eq:d3-ramification}
\end{equation}
then $\mathcal Z(P,Q,R)$ has an isolated point outside the diagonal.
\end{proposition}

\begin{proof}
Put $N=\deg R\geq1$.  Set $w=t^q$, $\widetilde w=T^q$, and write
$P=\overline P(x,w)$, $Q=\overline Q(x,w)$.  In each case, the map
\[
 \phi:\C^2_{x,w}\longrightarrow\C^2,\qquad
 (x,w)\longmapsto\bigl(\overline P(x,w),\overline Q(x,w)\bigr)
\]
is finite because its components $\overline P,\overline Q$, homogeneous
for weights $\deg x=1$ and $\deg w=q$, have only the origin as a common zero.
Indeed, on $P=0$, the form $Q$ equals
$xt^3$ in case $E$, and equals $2t^4$ or $-t^4$ in
case $H$.  The Jacobian of $\phi$ is
\[
 J_{x,w}(\overline P,\overline Q)=
 \begin{cases}-12w,&E,\\12w^2,&H.\end{cases}
\]

\smallskip
\noindent\emph{Components of $\mathcal Z(P,Q)$.}
In coordinates $[x:w:X:\widetilde w]$ on $\mathbb P(1,q,1,q)$,
let $D$ be the closure of the projective
classes represented by pairs of distinct points $(x,w),(X,\widetilde w)$
in a fiber of $\phi$.
The map
\[
 \rho_q:\mathbb P^3\longrightarrow\mathbb P(1,q,1,q),\qquad
 [x:t:X:T]\longmapsto[x:t^q:X:T^q]
\]
is well defined by the coordinate weights and finite because the
corresponding extension of homogeneous coordinate rings adjoins $q$-th
roots of $w,\widetilde w$.  We consider every curve component of
$\rho_q^{-1}(D)$.  If $R(x,t)-R(X,T)$ is
not identically zero on any component, its zeros on their union are isolated.

For $E$, the equations for $D$ are
\[
 w=\frac{(X-x)(X^2+2Xx+3x^2)}{12},\qquad
 \widetilde w=-\frac{(X-x)(3X^2+2Xx+x^2)}{12}.
\]
These equations define an isomorphism $\mathbb P^1_{[x:X]}\to D$,
whose inverse is the projection of $D$ to $[x:X]$, since $x,X$
cannot both vanish on $D$.  On $x=1$, put $\lambda=X$.  Then
\begin{equation}
 w=\frac{(\lambda-1)(\lambda^2+2\lambda+3)}{12},\qquad
 \widetilde w=-\frac{(\lambda-1)(3\lambda^2+2\lambda+1)}{12}.
                                                        \label{eq:d3-E-normalization}
\end{equation}
For $H$, on the chart $x=1$ with $X\neq0$, the fiber equation
$\overline P(1,w)=\overline P(X,\widetilde w)$ gives
\begin{equation}
 \widetilde w=\frac{1+3w-X^3}{3X}.                 \label{eq:d3-HW}
\end{equation}
Substitution in the fiber equation
$\overline Q(1,w)=\overline Q(X,\widetilde w)$ gives
$(X-1)\Gamma_H(X,w)/(9X^2)=0$, where
\begin{equation}
 \Gamma_H=18(X+1)w^2-12(X-1)(2X+1)w
 +(X-1)^2(X^3+3X^2+6X+2).                         \label{eq:d3-Gamma}
\end{equation}
After eliminating $\widetilde w$ by \eqref{eq:d3-HW}, the factor $X-1$
defines the diagonal, and $\Gamma_H=0$ defines $D$
in the $(X,w)$ chart.
Regarded as a quadratic in $w$ over $\C(X)$, $\Gamma_H$ has discriminant
$-72X^2(X-1)^2(X^2+4X+1)$, which is not a square in $\C(X)$
because $X^2+4X+1$ has two simple roots, neither equal to $0$ or $1$.
Hence $\Gamma_H$ is irreducible in $\C(X)[w]$.
Its coefficients have no common nonconstant factor in $\C[X]$,
so Gauss's lemma implies that it is irreducible in $\C[X,w]$ as well.
Thus $D$ is irreducible.

Let $C$ be the reduced union of the curve components of
$\rho_q^{-1}(D)$.  Each component maps onto $D$, since $\rho_q$
is finite and $D$ is irreducible.
The point $\rho_q(p_0)$ lies on $D$, and its inverse image is
$\{p_0\}$, so every component of $C$ passes through $p_0$.
Consequently, the fiber product decomposes as
\begin{equation}
 \bigl|\mathcal Z(P,Q)\bigr|
 =C\cup\bigcup_{\xi^q=1}L_\xi,
 \qquad L_1=\Delta.
 \label{eq:d3-fibre-components}
\end{equation}
No further curve component of $\mathcal Z(P,Q)$ is contained in $xX=0$.
Indeed, neither coordinate
pair $(x,t)$ nor $(X,T)$ is zero at a point of $\mathcal Z(P,Q)$,
and finiteness of $(P,Q)$
gives only finitely many points of $\mathcal Z(P,Q)$ over any fixed
$[x:t]$ or $[X:T]$.
The affine fiber product of the finite map $(P,Q)$ with itself is finite
over either copy of $\C^2$, so it has dimension two.
The ideal generated by its two defining differences
$P(x,t)-P(X,T)$ and $Q(x,t)-Q(X,T)$ therefore has height two in
$\C[x,t,X,T]$.
The projectivization $\mathcal Z(P,Q)$ of this affine fiber product
has pure dimension one.
Distinct lines $L_\xi$ meet only at $p_0$, and $C$ meets their
union only there.  Indeed, $\phi$ is locally biholomorphic for
$w\neq0$, so $D$ cannot meet the diagonal
$\{x=X,w=\widetilde w\}$ there.
On that diagonal, $w=0$ represents only $\rho_q(p_0)$.

\smallskip
\noindent\emph{Forms $R$ invariant under the cyclic action.}
Choose a primitive $q$-th root of unity $\xi$.  If the restriction of
$R(x,t)-R(X,T)$ to $L_\xi$ in \eqref{eq:d3-fibre-components} is
not identically zero and has a zero away from $p_0$, that zero is an isolated
collision.  First suppose that the restriction vanishes
identically.  Then $R=R_0(x,w)$ with $R_0\in\C[x,w]$ of weighted
degree $N$ for $\deg x=1$, $\deg w=q$.
The difference $R_0(x,w)-R_0(X,\widetilde w)$ is not identically zero
on $D$.  Otherwise $R_0$ would be constant on the fibers of $\phi$
over its regular values.  The finite extension
$\C(\overline P,\overline Q)\subset\C(x,w)$ is separable.
Constancy on these fibers would make all conjugates of $R_0$
over $\C(\overline P,\overline Q)$ coincide, so its minimal polynomial
over that field would have degree one and
$R_0\in\C(\overline P,\overline Q)$.
Finiteness of $\phi$ also makes $R_0$ integral over the polynomial
ring $\C[\overline P,\overline Q]$.  This ring is integrally
closed, so $R_0\in\C[\overline P,\overline Q]$.  Substituting
$w=t^q$ would give $R\in\C[P,Q]$, contrary to
\eqref{eq:d3-ramification}.
Since every component of $C$ maps onto $D$, none is contained
in $\Delta_R$.  Any zero of $R_0(x,w)-R_0(X,\widetilde w)$ on $D$
at a pair of distinct points gives, after taking $q$-th roots of
$w,\widetilde w$,
a point of $C\cap\Delta_R$ away from $p_0$.
By \eqref{eq:d3-fibre-components}, such a point is an isolated
collision.

In case $H$, weighted homogeneity of $R_0$ gives such a pair
as follows.  If $N$ is odd, every monomial of $R_0$ contains $x$.
Thus $(0,w)$ and $(0,-w)$ have the same values of
$\overline P,\overline Q,R_0$ for $w\neq0$.
If $N$ is even, the exponent of $x$ in every monomial of $R_0$ is even.
The points $(x,-x^2/3)$ and $(-x,-x^2/3)$ then have the same
three values for $x\neq0$.

In case $E$, use \eqref{eq:d3-E-normalization} to define the polynomial
\[
 F(\lambda)=R_0(1,w(\lambda))
             -R_0(\lambda,\widetilde w(\lambda)).
\]
The polynomial $F$ is nonzero and satisfies $F(1)=0$.
The chain rule and \eqref{eq:d3-ramification} give
\[
 w\mid J_{x,w}(\overline P,R_0),\qquad
 w\mid J_{x,w}(\overline Q,R_0).
\]
Define the polynomials
\[
 \Lambda_P=\frac{J_{x,w}(\overline P,R_0)}w,\qquad
 \Lambda_Q=\frac{J_{x,w}(\overline Q,R_0)}w.
\]
Since $J_{x,w}(\overline P,\overline Q)=-12w$, Cramer's rule gives
\[
 dR_0=\frac{\Lambda_Q}{12}\,d\overline P
       -\frac{\Lambda_P}{12}\,d\overline Q.
\]
For a polynomial $\Phi\in\C[x,w]$, write
$[\Phi]=\Phi(1,w(\lambda))
-\Phi(\lambda,\widetilde w(\lambda))$.
On the chart $x=1$ of $D$, the common values of $\overline P$ and
$\overline Q$ are $1+3w(\lambda)$ and $1+4w(\lambda)$.
Pulling back the identity for $dR_0$ along the maps from $\C$ to
$\C^2$ given by $\lambda\mapsto(1,w(\lambda))$ and
$\lambda\mapsto(\lambda,\widetilde w(\lambda))$, and subtracting, yields
\begin{equation}
 F'(\lambda)=\frac{w'(\lambda)}{12}
                 \bigl(3[\Lambda_Q]-4[\Lambda_P]\bigr).
 \label{eq:d3-derivative}
\end{equation}
By \eqref{eq:d3-derivative}, $w'(\lambda)=(3\lambda^2+2\lambda+1)/12$ divides
$F'(\lambda)$, and the roots of $w'(\lambda)$ are different from $1$.
If $1$ were the only zero of $F$, then
$F=c(\lambda-1)^k$ with $c\neq0$ and $k\geq1$, contradicting
this divisibility.  Hence $F$ has a zero with $\lambda\neq1$,
which gives a pair of distinct points with equal values of
$\overline P,\overline Q,R_0$.

\smallskip
\noindent\emph{The remaining forms $R$.}
It remains to consider a nonzero restriction of $R(x,t)-R(X,T)$ to $L_\xi$ whose
only zero is $p_0$.  This restriction is a multiple of $t^N$.
Comparison of powers of $t$ then gives
\[
 R(x,t)=R_0(x,t^q)+ct^N,
 \qquad c\neq0,\qquad q\nmid N.
\]
Here $R_0\in\C[x,w]$.
In case $E$, \eqref{eq:d3-E-normalization} gives one smooth branch
of $D$ at $\rho_q(p_0)$, with $w$ a local parameter and
$\widetilde w/w\to-1$.
In case $H$, the tangent cone of the curve $\Gamma_H=0$, with
$\Gamma_H$ defined in \eqref{eq:d3-Gamma}, at $(X,w)=(1,0)$
is defined by the quadratic form
\[
 12\bigl((X-1)^2-3(X-1)w+3w^2\bigr).
\]
The two distinct lines of this tangent cone give two smooth branches of $D$,
$X=1+\lambda w+O(w^2)$, where $\lambda^2-3\lambda+3=0$.
By \eqref{eq:d3-HW}, the corresponding limits of
$\widetilde w/w$ are $1-\lambda$.
Thus $w$ is a local parameter on every branch of $D$ at $\rho_q(p_0)$,
and all these limits are nonzero and different from $1$.
Taking $q$-th roots of $w,\widetilde w$ gives, on every local branch
of $C$ at $p_0$,
\[
 w=t^q,\qquad T=tU(t^q),\qquad U\in\C\{w\}^*.
\]

We compare $R(x,t)-R(X,T)$ with $t^N-T^N$.
The latter vanishes at the following points of $C$:
\begin{equation}
\begin{array}{c|c|c|c}
 &X&t=T&x\\ \hline
 E&\omega x&1&x^3=-4,\quad \omega^3=1,\ \omega\neq1\\
 H&-x&1&x^2=-3.
\end{array}
\label{eq:d3-monomial-pairs}
\end{equation}
For $E$, the values of $P$ agree and both values of $Q$ vanish.
For $H$, both values of $P$ vanish and the values of $Q$ agree.
At the images of these points under $\rho_q$, the coordinate pairs
$(x,w)$ and $(X,\widetilde w)$ are distinct, so the points lie on
$C$ away from all the lines $L_\xi$.
The rational function $\widetilde w/w$ on $D$ equals $1$ at
their images, whereas its limits at $\rho_q(p_0)$ differ from $1$.
It is therefore nonconstant.
If $t^N-T^N$ vanished on a component of $C$, then $T/t$ would
be a constant $N$th root of unity there.  Since that component
maps onto $D$, the function $\widetilde w/w=(T/t)^q$ would
be constant on $D$, a contradiction.
Thus no component of $C$ is contained in $\Delta_{t^N}$.

On each local branch of $C$ at $p_0$, in the chart $x=1$, the
difference $R(x,t)-R(X,T)$ is
\[
 A(t^q)+ct^N\bigl(1-U(t^q)^N\bigr),
 \qquad A\in\C\{w\}.
\]
Here $A(w)$ is the restriction of $R_0(x,w)-R_0(X,\widetilde w)$
to the corresponding local branch of $D$.
The exponents in the two summands have different residues modulo $q$,
so their leading terms cannot cancel.  Hence
\[
 \ord_t\bigl(R(x,t)-R(X,T)\bigr)
 \leq\ord_t(t^N-T^N).
\]
Since every component of $C$ passes through $p_0$, the inequality
implies that none is contained in $\Delta_R$.
For a hypersurface $\mathcal H$ containing no component of $C$, we write
$i_p(C,\mathcal H)$ for the sum of the vanishing orders of a local equation
of $\mathcal H$ along the normalized branches of $C$ over $p$.
We write $\deg C$ for the number of intersections of $C$
with a general hyperplane.  Both differences have degree $N$.
The curve--hypersurface form of B\'ezout's theorem, together
with the zero in \eqref{eq:d3-monomial-pairs} away from $p_0$, gives
\[
 i_{p_0}(C,\Delta_R)
 \leq i_{p_0}(C,\Delta_{t^N})<N\deg C.
\]
B\'ezout's theorem gives total intersection number $N\deg C$
for $C$ and $\Delta_R$, so the strict inequality leaves an
intersection point away from $p_0$.  Since no component of $C$
is contained in $\Delta_R$, this point is isolated in
$C\cap\Delta_R$.  By \eqref{eq:d3-fibre-components}, $C$ meets
the diagonal and all remaining components of $\mathcal Z(P,Q)$
only at $p_0$.  The point is therefore outside the diagonal and
isolated in $\mathcal Z(P,Q,R)$.
\end{proof}

\subsection{Exclusion of order three}

The following lemma prepares the third coordinate when
$m_{\boldsymbol n}=3$ and when $m_{\boldsymbol n}=2$.
Once the initial forms of the first two coordinates define a finite map,
we subtract polynomials
in the first two coordinates from the third coordinate, as in
Step~\ref{step:third-coordinate}.  On a fiber of the first two
coordinates, each subtracted function is constant, so every difference
of values of the third coordinate is preserved.  When the map given by
the first two coordinates has degree greater than one, the lemma shows
that injectivity forces this process to reach an initial form of the third
coordinate outside the polynomial algebra generated by the first two
initial forms.

For fixed weights $(w_x,w_y)$, let $G,U\in\mathfrak m$ have initial
forms $P,Q$ with no common zero except the origin, and put $A=\C\{G,U\}$.
Give $A$ the filtration induced by weighted order on $B$, and
write $\operatorname{gr}A$ for its \emph{associated graded ring}.
We identify it with the subalgebra of $\C[x,y]$ generated by the
initial forms of elements of $A$.
Let $H\in B$.

\begin{lemma}
\label{lem:finite-plane-degree}
\label{lem:filtered-escape}
The map germ $(G,U):(\C^2,0)\to(\C^2,0)$ is finite, and its fibers
over regular values consist of $i_0(G,U)$ distinct points.  Moreover,
$\operatorname{gr}A=\C[P,Q]$.
If $i_0(G,U)>1$ and $H\in A$, then $(G,U,H)$ is not
injective.  If $H\notin A$, successive subtraction from $H$ of polynomials
in $G,U$ that cancel the initial form of the current remainder terminates at a
remainder whose initial form lies outside $\C[P,Q]$.
\end{lemma}

\begin{proof}
Suppose that nonzero common zeros of $G,U$ approach the origin.
For each such zero $(x,y)$, set $t=\max\{|x|^{1/w_x},|y|^{1/w_y}\}$.
Rescaling to $(x/t^{w_x},y/t^{w_y})$ and taking a subsequence would
give a nonzero common zero of $P,Q$.  Thus $(G,U)$ is
finite by \cite[Theorem~I.1.66]{GreuelLossenShustin}.
The substitution homomorphism $\C\{S,T\}\to B$, $F\mapsto F(G,U)$,
makes $B$ a finitely generated module over its image $A$, so $\dim A=\dim B=2$.
Its kernel is a prime ideal of height zero in the domain
$\C\{S,T\}$, and is therefore zero.
Thus $A\cong\C\{S,T\}$.  Both $A$ and $B$ are regular local rings
of dimension two, and $G,U$ form a $B$-regular sequence.
Hence $B$ is free over
$A$ by \cite[Theorem~B.8.11 and Proposition~B.3.5]{GreuelLossenShustin},
of rank $\dim_\C B/(G,U)=i_0(G,U)$.  In characteristic zero
the fibers of $(G,U)$ over regular values have this many distinct points
\cite[Theorem~I.1.81(2)]{GreuelLossenShustin}.
A function in $A$ is constant on these fibers, proving that $(G,U,H)$
is not injective when $i_0(G,U)>1$ and $H\in A$.

The finite polynomial map $(P,Q)$ makes $\C[x,y]$ integral over
$\C[P,Q]$.  Hence $\C[P,Q]$ has dimension two, and its two
generators $P,Q$ are algebraically independent.
Give $S,T$ weights equal to the weighted orders of $G,U$,
respectively.  For a nonzero $F\in\C\{S,T\}$, let $F_d$
be its initial form for these weights.  Algebraic independence gives
$F_d(P,Q)\neq0$, so the initial form of $F(G,U)$ is precisely
$F_d(P,Q)$.  This proves $\operatorname{gr}A=\C[P,Q]$.
Each cancellation strictly increases the weighted order of the remainder.
Suppose that the process did not terminate.  With
$\mathfrak m_A=(G,U)A$, the ideal $\mathfrak m_A B=(G,U)B$
is $\mathfrak m$-primary, so
$\mathfrak m^c\subset\mathfrak m_A B$ for some $c\geq1$.
Since the weighted orders of the remainders are unbounded,
$H\in A+\mathfrak m^{cn}\subset A+\mathfrak m_A^nB$ for every
$n$.  The class of $H$ in the finite $A$-module $B/A$
then vanishes by the Krull intersection theorem
\cite[Theorem~B.4.2]{GreuelLossenShustin}, contradicting $H\notin A$.
\end{proof}

Applied to $v,u$ with weights $(a,1)$, \Cref{lem:filtered-escape} supplies
an initial form of the third coordinate of $\boldsymbol n$ outside the
algebra generated by the initial forms of $v,u$.
If this form is a polynomial in $x,y^a$,
\Cref{prop:d3-cubic-collision} applies after the substitution $t=y^a$.
Otherwise, multiplying $y$ by an $a$-th root of unity or choosing the
points in \eqref{eq:d3-monomial-pairs} gives a collision.

\begin{theorem}
\label{thm:dthree}
Under the hypotheses of \Cref{sec:input}, one has
$m_{\boldsymbol n}\neq3$.
\end{theorem}

\begin{proof}
Assume $m_{\boldsymbol n}=3$, and use the coordinates of
\Cref{prop:d3-models}.  Put
$\widehat P=P(x,y^a)$, $\widehat Q=Q(x,y^a)$, and denote the
third coordinate by $w$, so $\boldsymbol n=(v,u,w)$.
The map $(\widehat P,\widehat Q)$ is finite.  By \Cref{lem:finite-plane-degree},
$(v,u)$ has degree $m_X=12a>1$.  Injectivity of $\boldsymbol n$ therefore implies
$w\notin\C\{v,u\}$.  By \Cref{lem:filtered-escape}, subtracting a
polynomial in $v,u$ from $w$ gives an initial form
\[
 \mathcal R=\operatorname{in}_{(a,1)}w
 \notin\C[\widehat P,\widehat Q].
\]
Write the weight of $\mathcal R$ as $aN+r$, with $0\leq r<a$.  Then
$\mathcal R=y^rS(x,y^a)$, where $S(x,t)$ is homogeneous of
degree $N$.

If $r=0$, taking initial forms with weights $(a,1)$ in
\eqref{eq:d3-minors} and using
\[
 J_{x,y}\bigl(P(x,y^a),S(x,y^a)\bigr)
 =ay^{a-1}J_{x,t}(P,S)(x,y^a)
\]
gives $t^5\mid J_{x,t}(P,S),J_{x,t}(Q,S)$, including when either
Jacobian is zero.  Since $S\notin\C[P,Q]$,
\Cref{prop:d3-cubic-collision} gives an isolated non-diagonal point
of $\mathcal Z(P,Q,S)$.  The map
\[
 \rho_a:\mathbb P(a,1,a,1)\longrightarrow\mathbb P^3,
 \qquad [x:y:X:Y]\longmapsto[x:y^a:X:Y^a]
\]
is well defined, since its coordinate polynomials have weighted degree $a$.
It is finite because the corresponding extension of homogeneous
coordinate rings adjoins $a$-th roots of $t,T$, and surjective because
these roots exist over $\C$.  Substituting $t=y^a$
and $T=Y^a$ in the equations defining $\mathcal Z(P,Q,S)$ gives
$\mathcal Z(\widehat P,\widehat Q,\mathcal R)
=\rho_a^{-1}(\mathcal Z(P,Q,S))$.  The inverse image under $\rho_a$ of the
isolated collision is therefore nonempty and consists of isolated points.
These points lie outside the diagonal in $\mathbb P(a,1,a,1)$, since
$\rho_a$ maps this diagonal into the diagonal in $\mathbb P^3$.

Suppose $0<r<a$, and choose a primitive $a$-th root of unity $\zeta$.
On $(X,Y)=(x,\zeta y)$, the differences of the values of $\widehat P$
and $\widehat Q$ vanish, and $\mathcal R(x,y)-\mathcal R(X,Y)$ is
$(1-\zeta^r)y^rS(x,y^a)$.  If $S$ has a zero
$[x:t]\in\mathbb P^1$ with $t\neq0$, that zero gives a collision of
$(\widehat P,\widehat Q,\mathcal R)$ with
$y\neq0$.  It is isolated: by \eqref{eq:d3-PQjac}, the differences
for $\widehat P,\widehat Q$ locally define the curve $(X,Y)=(x,\zeta y)$.
On this curve in the chart $y=1$, the difference for $\mathcal R$
restricts to a nonzero multiple of $S(x,1)$.

If $S$ has no zero $[x:t]\in\mathbb P^1$ with $t\neq0$, the nonzero polynomial
$S(x,1)\in\C[x]$ has no zero and is therefore constant.
Homogeneity of degree $N$ then gives $S=ct^N$, with $c\neq0$,
allowing $N=0$.  In the chart $y=1$, take $x,X$ from
\eqref{eq:d3-monomial-pairs} and put $Y=1$.
The initial forms $\widehat P,\widehat Q$ agree at these distinct points by the
calculation following \eqref{eq:d3-monomial-pairs}.  The values of
$\mathcal R=cy^{aN+r}$ also agree because $Y=y$.
At these points, the derivative of $\mathcal R(x,y)-\mathcal R(X,Y)$
with respect to $Y$ is $-c(aN+r)\neq0$, and the Jacobian determinant of
the differences for $\widehat P,\widehat Q$ with respect to $x,X$ is
\[
 36x^2(1-\omega^2)\neq0\quad(E),\qquad48x\neq0\quad(H).
\]
Thus the common zero is isolated.

In every case, $\mathcal Z(\widehat P,\widehat Q,\mathcal R)$
has an isolated non-diagonal point.  The positive integral weights
$(a,1)$ allow \Cref{lem:weighted-lifting} to produce collisions
for $\boldsymbol n$, contradicting injectivity.
\end{proof}
\section{Exclusion of
\texorpdfstring{$m_{\boldsymbol{n}}=2$}{m(n)=2}}
\label{sec:dtwo}

Assume $m_{\boldsymbol n}=2$ and choose the component $\Gamma_i$
and coordinates from \Cref{prop:common-preparation}.  Then
\[
 \Gamma_i=\{y=0\},\qquad a_i=2,\qquad
 v=E(x^2+c_0(y)),\qquad \ord_yc_0\geq2A_2.
\]
We first determine the quadratic terms of $\boldsymbol n$ in
\Cref{prop:d2-quadratic-jet}.  These terms will allow an analytic change
of coordinates that preserves $\Gamma_i=\{y=0\}$ and puts $v$ in the
form stated in \Cref{lem:d2-morse}.
Divisibility of the Jacobian minors of $\boldsymbol n$ then reduces this map to the
form in \Cref{prop:d2-cyclic-preparation}, which
\Cref{prop:d2-cyclic-separation} excludes.
We apply Steps~\ref{step:root-separation} and
\ref{step:third-coordinate} successively to refine the
coordinates.  In the final form of \Cref{prop:d2-cyclic-preparation},
the third coordinate must take different values at distinct points in
each fiber of the finite map $(P,Q)$ formed by the first two coordinates.
We find a pair violating this condition by studying a local component
of the fiber product, with a second choice of weights when necessary.
Put
\[
 A_j=\#\{p\in\mathcal C:m_{\boldsymbol n,p}=j\}\quad(j=1,2),
 \qquad \mu_n=\{\zeta\in\C^*:\zeta^n=1\}\quad(n\geq1).
\]
For every $p\in\mathcal C$, the multiplicity $m_{\boldsymbol n,p}$
is one or two, and the origin is counted by $A_2$.  Equations \eqref{eq:input-point-basis} and
\eqref{eq:input-polar} become
\begin{align}
 m_X&=4A_2+A_1,& M_{\boldsymbol n}&=2A_2+A_1,&
 \delta(V,0)&=A_2,                                  \label{eq:d2-counts}\\
 \sum_i a_i(m_i-1)&\leq6A_2+A_1-b,&&&                 \label{eq:d2-polar}
\end{align}
where $b\in\{1,2\}$.

\subsection{The quadratic jet}

Let
\[
 \mathcal J_2=\operatorname{span}_{\C}
 \{(n_1)_2,(n_2)_2,(n_3)_2\},
\]
where $(F)_2$ is the \emph{homogeneous quadratic part} of $F$.
Since $(d\boldsymbol n)_0=0$, the triple
$((n_1)_2,(n_2)_2,(n_3)_2)$ is the \emph{quadratic jet}
of $\boldsymbol n$.
To show that $\mathcal J_2$ is generated by the square of a linear form,
we first prove the following two lemmas.  The first will also exclude the second numerical
alternative in \Cref{prop:d2-second-alternative}.  Put
\begin{equation}
 \Phi(s,\lambda)=s^2-\lambda^2,\qquad
 \Psi(s,\lambda)=2s^3-3s\lambda^2.
 \label{eq:d2-cubic-pair}
\end{equation}

\begin{lemma}
\label{lem:d2-cubic-orbit}
If $R\in\C[s,\lambda]$ is homogeneous of degree $r\geq4$ and
$R\notin\C[\Phi,\Psi]$, then
$\mathcal Z(\Phi,\Psi,R)$ has an isolated point outside the diagonal.
\end{lemma}

\begin{proof}
The identities $s^3-3\Phi s+\Psi=0$ and
$\lambda^2=s^2-\Phi$ make the polynomial map
$(\Phi,\Psi):\C^2\to\C^2$ finite.
In $\mathbb P^3$, with coordinates $[s:\lambda:s':\lambda']$, set
$D_\pm=\{s'=s,\lambda'=\pm\lambda\}$.  On the locus
$\Phi(s,\lambda)=\Phi(s',\lambda')$, factoring
$\Psi(s,\lambda)-\Psi(s',\lambda')$ gives
\[
 \mathcal Z(\Phi,\Psi)=D_+\cup D_-\cup\Gamma,
\]
as sets, where $\Gamma$ is defined by the equations
\begin{equation}
 (s')^2+ss'-2s^2+3\lambda^2=0,\qquad
 (\lambda')^2=(s')^2-s^2+\lambda^2.                 \label{eq:d2-Gamma}
\end{equation}
The normalization map of $\Gamma$ is the following morphism from the
smooth conic
$\mathcal N=\{\rho^2=3(\alpha^2+\beta^2)\}\subset\mathbb P^2$:
\[
 \begin{aligned}
 \mathcal N&\longrightarrow\Gamma,\\
 [\rho:\alpha:\beta]&\longmapsto
 [-(\alpha^2+2\beta^2):\rho\beta:
   2\alpha^2+\beta^2:\rho\alpha].
 \end{aligned}
\]
Indeed, the four quadratic forms defining this morphism have no common
zero on $\mathcal N$, and their values satisfy \eqref{eq:d2-Gamma}.
The morphism has inverse $[s'-s:\lambda':\lambda]$ away from
$o=[1:0:1:0]$.  The point $o$ is the only point of $\Gamma$ with
$s'=s$, and its preimages are $p_\pm=[0:\pm i:1]$.
In particular, $\Gamma\cap(D_+\cup D_-)=\{o\}$.

For $d\geq0$, write $\mathcal O_{\mathcal N}(d)$ for the restriction
of the line bundle $\mathcal O_{\mathbb P^2}(d)$ to $\mathcal N$.
A homogeneous form of degree $d$ restricts to a section of this line
bundle.
The pullback $\mathcal R$ of $R(s,\lambda)-R(s',\lambda')$ under the
normalization map $\mathcal N\to\Gamma$ is a nonzero section of
$\mathcal O_{\mathcal N}(2r)$.
If it vanished identically, $R$ would agree on every ordered pair in
a general fiber of $(\Phi,\Psi)$ having different $s$-coordinates.  Such a fiber has
three distinct $s$-coordinates and two points above each, so $R$ would be
constant on the whole fiber.  Constancy on general fibers and
separability of $\C(\Phi,\Psi)\subset\C(s,\lambda)$ would give
$R\in\C(\Phi,\Psi)$.  Finiteness makes $R$ integral over
$\C[\Phi,\Psi]$, which is a polynomial ring and hence integrally
closed.  Thus $R\in\C[\Phi,\Psi]$, contrary to the hypothesis.

The zero divisor of a nonzero section records its zeros with their
orders of vanishing.
Suppose the zeros of $\mathcal R$ lie only at $p_\pm$.  Their orders sum to $4r$
by B\'ezout's theorem for the conic and a form of degree $2r$
\cite[Section~2.6, p.~54]{CasasAlvero}.
Interchanging $(s,\lambda)$ and $(s',\lambda')$ exchanges $p_+$ and $p_-$, so both
orders are $2r$.  Since the pullback of $s'-s$ is $\rho^2$, which has order two at both
points, both $\mathcal R$ and the pullback of $(s'-s)^r$ are
sections of $\mathcal O_{\mathcal N}(2r)$ with zero divisor
$2r(p_++p_-)$.  Their quotient is a nowhere vanishing regular
function on the projective conic $\mathcal N$, hence a
nonzero constant.
Thus $\mathcal R=c(s'-s)^r$, with $c\neq0$,
where the right-hand side denotes the pullback to $\mathcal N$.
Interchanging the two factors negates $\mathcal R$ and multiplies
$(s'-s)^r$ by $(-1)^r$, so $r$ is odd.
In the fiber $(\Phi,\Psi)=(1/3,0)$, take one
point above each of $s=-1,0,1$.  Summing the identity
$R(s,\lambda)-R(s',\lambda')=c(s'-s)^r$ for the three
ordered pairs with $s$-coordinates $(-1,0)$, $(0,1)$, and $(1,-1)$
yields $0=c(1+1-2^r)$, a contradiction.
Thus $\mathcal R$ has a zero away from $p_\pm$.  Its image is an isolated zero
of $R(s,\lambda)-R(s',\lambda')$ on $\Gamma$, disjoint from the other components of
$\mathcal Z(\Phi,\Psi)$, and outside the
diagonal $D_+$.
\end{proof}

The second lemma concerns plane sections.  For the choice of
representatives of deformations in a Milnor ball, see
\cite[Definition~2.2]{SteenbrinkSemicontinuity}.

\begin{lemma}
\label{lem:d2-forbidden-splitting}
Let $(W,0)$ be a reduced plane curve germ of multiplicity four with
two branches.  There is no one-parameter flat deformation of $(W,0)$
represented by plane curves $W_s$ in a fixed Milnor ball such that
$W_0$ is singular only at the origin, all singular points of $W_s$
tend to the origin as $s\to0$, and the general fiber $W_s$, for
sufficiently small $s\neq0$, has only unibranch singularities of
multiplicity two and satisfies
\begin{equation}
 \delta(W,0)=1+\sum_{Q\in\operatorname{Sing}(W_s)}\delta(W_s,Q).
 \label{eq:d2-forbidden-delta}
\end{equation}
\end{lemma}

\begin{proof}
We use the spectrum $\operatorname{Sp}$ as defined in
\cite[(2.1)]{SteenbrinkSemicontinuity}, with spectral numbers of
plane curve singularities in $(-1,1)$ counted with multiplicity.
Let $F=0$ be a reduced equation for $(W,0)$ in local coordinates $(x,y)$.
Write $F=F_4+H$,
where $F_4$ is homogeneous of degree four and
$H\in(x,y)^5\C\{x,y\}$.
Put $G_t=F_4+t(x^4+y^4)$.
Its discriminant is a nonzero polynomial in $t$, with leading
coefficient the nonzero discriminant of $x^4+y^4$.
Thus $G_t$ is squarefree for all sufficiently small $t\neq0$.
For each such $t$, $G_t$ has an isolated critical point at the origin,
while $uH$ has order at least five for every $u\in\C$.  Hence
\cite[Corollary~I.2.18]{GreuelLossenShustin} gives
$\mu(G_t+uH)=\mu(G_t)=9$.
Invariance of the spectrum under deformations with constant Milnor
number \cite[Theorem~2.8]{SteenbrinkSemicontinuity}
therefore identifies $\operatorname{Sp}(F+t(x^4+y^4))$ with
$\operatorname{Sp}(G_t)$.  The spectrum formula for the homogeneous polynomial $G_t$
\cite[p.~165]{SteenbrinkSpectrum} gives
\[
 \operatorname{Sp}(F+t(x^4+y^4))
 =\left\{\frac{i+j}{4}-1:1\leq i,j\leq3\right\}.
\]
Lower semicontinuity of the smallest spectral number
\cite[Theorem~2.11]{SteenbrinkSemicontinuity}, applied to
the deformation $F+t(x^4+y^4)$, yields
$\min\operatorname{Sp}(W)\leq-1/2$.
Spectral symmetry \cite[p.~167]{SteenbrinkSpectrum} gives
$-\min\operatorname{Sp}(W)\geq1/2$ as another spectral number.
These two distinct spectral numbers lie outside $(-1/2,1/2)$,
so at most $\mu(W,0)-2$ spectral numbers lie in that interval.

By \cite[Theorem~I.2.48]{GreuelLossenShustin}, each germ
$(W_s,Q)$, with $Q\in\operatorname{Sing}(W_s)$, has an
equation $z^2-w^{q_Q}=0$ in suitable local coordinates, with
$q_Q\geq3$ odd by irreducibility.
Its spectrum is
\[
 \left\{-\frac12+\frac j{q_Q}:1\leq j<q_Q\right\}
 \subset(-1/2,1/2)
 \qquad\text{\cite[p.~165]{SteenbrinkSpectrum}}.
\]
Under the stated assumptions on the representatives $W_s$,
spectral semicontinuity applies to the number of spectral numbers
in each open interval of length one
\cite[Remark~2.10]{SteenbrinkSemicontinuity}.  Hence
\[
 \sum_Q\mu(W_s,Q)
 \leq\#\bigl(\operatorname{Sp}(W)\cap(-1/2,1/2)\bigr)
 \leq\mu(W,0)-2.
\]
Applying $\mu=2\delta-b+1$ to $(W,0)$ and each $(W_s,Q)$, with
$b$ denoting the number of branches of the germ in question, and using
\eqref{eq:d2-forbidden-delta} instead gives
$\sum_Q\mu(W_s,Q)=2\delta(W,0)-2=\mu(W,0)-1$, a contradiction.
\end{proof}

We will also use the \emph{splitting lemma with one coordinate fixed}.
For a holomorphic function germ $f\in\C\{x,y\}$ with
$f_x(0)=0$ and $f_{xx}(0)\neq0$, the implicit function theorem gives
a holomorphic germ $\xi(y)$ with $\xi(0)=0$ and $f_x(\xi(y),y)=0$.
Taylor expansion gives
\[
 f(x,y)=f(\xi(y),y)+(x-\xi(y))^2U(x,y),\qquad U(0)\neq0.
\]
Thus the change of coordinates $(x,y)\mapsto(X,y)$, where
$X=(x-\xi(y))\sqrt{U(x,y)}$, puts $f$ in the form $X^2+h(y)$,
where $h(y)=f(\xi(y),y)$, while preserving $y$.  This proves the form
of \cite[Theorem~I.2.47]{GreuelLossenShustin} that preserves the parameter $y$, as needed below.

The preceding lemmas now determine $\mathcal J_2$.

\begin{proposition}
\label{prop:d2-quadratic-jet}
There is a nonzero linear form $\ell$ such that
$\mathcal J_2=\operatorname{span}_{\C}\{\ell^2\}$.
\end{proposition}

\begin{proof}
If $\dim\mathcal J_2=3$, choose target coordinates so that the quadratic
initial forms of the components of $\boldsymbol n$ are $x^2,xy,y^2$.
Their Jacobian minors $2x^2,4xy,2y^2$ have no common factor, contradicting the
nontrivial common divisor in \eqref{eq:input-divisor}.

Suppose $\dim\mathcal J_2=2$ and all elements of $\mathcal J_2$ have a nonconstant common factor,
so $\mathcal J_2=x\operatorname{span}_{\C}\{x,y\}$ in suitable coordinates.
By \Cref{prop:common-preparation}, the chosen component $\Gamma_i$
is smooth and $i_0(v,\Gamma_i)=2$.  Let $L$ be the linear initial
form of a reduced defining function of $\Gamma_i$.  Divisibility of the
Jacobian minors gives
\[
 L^{m_{\mathrm{tr}}-1}\mid J(x^2,xy)=2x^2.
\]
Thus $L$ is proportional to $x$.  Since the quadratic part of $v$
is also divisible by $x$, the restriction of $v$ to $\Gamma_i$ has order
greater than two, a contradiction.

Suppose instead that the two generators of $\mathcal J_2$ have no common zero in $\mathbb P^1$.
The ideal $I$ has no base point after the blowup at the origin, so
$A_2=1,A_1=0,m_X=4,b=2$.  Then $M_{\boldsymbol n}=2$, and
\Cref{prop:common-preparation} gives $m_{\mathrm{tr}}=2$.
Hence every $m_i=2$.  Write $q_v,q_u$ for the coprime quadratic
initial forms of $v,u$.

The two branches of $V$ are smooth.  On each, $u$ has order two
because $q_u$ does not vanish at either root of $q_v$.  The images of the two branches of $V$ under $\boldsymbol n$
are the branches of $W$ and therefore have multiplicity two, so $W$
has multiplicity four and two branches.

The sections $W_s=X\cap\{\lambda_1=s\}$ form a one-parameter flat
deformation of $(W,0)$ because the local ring of the irreducible surface
$X$ at the origin is torsion-free over the discrete valuation ring
$\C\{\lambda_1\}$ \cite[Theorem~B.8.11]{GreuelLossenShustin}.
By \cite[Proposition~2.10 and Section~2.2]{PortillaReduction},
the germs $(W_s,Q)$ at the nearby singular points are the transverse
branches of $X$ along $\Sigma$, and
\[
 \delta(W,0)=\delta(V,0)+
 \sum_{Q\in\operatorname{Sing}(W_s)}\delta(W_s,Q).
\]
Choose representatives of $W_s$ transverse to the boundary of a Milnor ball for $(W,0)$.
All nearby singular points tend to the origin since
$\Sigma\cap\{\lambda_1=0\}=\{0\}$.  Since $\delta(V,0)=1$ and
every $m_i=2$, this contradicts \Cref{lem:d2-forbidden-splitting}.

It remains to exclude a one-dimensional $\mathcal J_2$ with a
squarefree generator.  The two simple zeros of this generator in $\mathbb P^1$ yield base
points $p$ with $m_{\boldsymbol n,p}=1$, so $A_2=1,A_1\geq2,b=2$.  For the chosen component $\Gamma_i$ of the critical locus
with $m_i=m_{\mathrm{tr}}$,
\eqref{eq:input-knot} and \eqref{eq:d2-polar} give
\[
 2(A_1+1)\leq a_i(m_{\mathrm{tr}}-1)\leq A_1+4.
\]
The inequalities $2(A_1+1)\leq A_1+4$ and $A_1\geq2$ give
$A_1=2$.  \Cref{prop:common-preparation} gives $a_i=2$ and a smooth
$\Gamma_i$.  Since $m_{\mathrm{tr}}\geq M_{\boldsymbol n}=4$,
the displayed inequalities force $m_{\mathrm{tr}}=4$.
Take a reduced defining function of $\Gamma_i$ as a coordinate $\lambda$.  The splitting lemma
with $\lambda$ fixed, followed by multiplication of $\lambda$ by a unit,
gives $v=s^2-\lambda^2=\Phi$, since the restriction of $v$ to $\lambda=0$
has order two and the Hessian of $v$ is nondegenerate.

After subtracting scalar multiples of $v$, the other two coordinate functions of $\boldsymbol n$ have order
at least three.  Both branches of $W$ then have multiplicity at
least three.  If both remaining coordinates had order at least
four, each image branch would have multiplicity at least four,
contrary to $\mult_0W=m_X=6$.
Thus some remaining coordinate has a nonzero cubic part $H_3$.  The divisibility
$\lambda^3\mid J_{s,\lambda}(\Phi,H_3)$ forces
$H_3\in\operatorname{span}_{\C}\{2s^3-3s\lambda^2\}$, by comparing the coefficients of $\lambda^0,\lambda^1,\lambda^2$
in $J_{s,\lambda}(\Phi,H_3)$.
After a linear change of target coordinates, the homogeneous parts of
least degree in the first two coordinate functions of $\boldsymbol n$
are $\Phi,\Psi$, and the third coordinate has order at least four.
The initial forms $\Phi,\Psi$ have degrees two and three and no
common zero in $\mathbb P^1$, so their intersection multiplicity is six.
By injectivity and \Cref{lem:filtered-escape}, subtracting a polynomial
in the first two coordinate functions from the third gives a third
coordinate whose initial form
$R\notin\C[\Phi,\Psi]$ has degree at least four.
\Cref{lem:d2-cubic-orbit,lem:weighted-lifting} contradict injectivity.
Thus $\mathcal J_2$ is spanned by the square of a linear form.
\end{proof}

\subsection{Divisibility of the Jacobian minors and the numerical alternatives}

We now refine the coordinates of \Cref{prop:common-preparation} while keeping $\Gamma_i$
as the axis $y=0$.  Since $\mathcal J_2$ is generated by the square of
a linear form, the splitting lemma applies to $v$ with $y$ fixed.  Divisibility of the
Jacobian minors along $\Gamma_i$ will then give a recurrence for
the coefficients of the coordinate functions of $\boldsymbol n$.  Put $k=2A_2-b+2$.

\begin{lemma}
\label{lem:d2-morse}
The integer $k$ satisfies $k\geq3$.  There are local coordinates $x,y$
such that
\[
 v=x^2+y^k,\qquad I(\Gamma_i)=(y).
\]
Here $I(\Gamma_i)$ is the ideal of functions in $B$ vanishing on $\Gamma_i$.
\end{lemma}

\begin{proof}
Return to the coordinates of \Cref{prop:common-preparation}.
Since $v(x,0)$ has order two, $v_{xx}(0)\neq0$.
The splitting lemma with $y$ fixed gives $v=x^2+h(y)$ and
preserves $\Gamma_i=\{y=0\}$.  By \Cref{prop:d2-quadratic-jet},
the Hessian of $v$ has rank one, so $\ord_y h\geq3$.
Reducedness of $v$ gives $h\neq0$, since $h=0$ would give
$v=x^2$.  The Jacobian ideal of $v$ is $(x,h'(y))\subset\C\{x,y\}$,
so $\mu(V,0)=\dim_\C\C\{y\}/(h'(y))=\ord_yh-1$, and
\eqref{eq:d2-counts} and \eqref{eq:input-milnor} yield
\[
 \ord_yh=\mu(V,0)+1=2A_2-b+2=k.
\]
Since $h(y)/y^k$ is a unit, replacing $y$ by
$y(h(y)/y^k)^{1/k}$ gives the required form and preserves
$\Gamma_i=\{y=0\}$.
\end{proof}

Fix the coordinates and component of \Cref{lem:d2-morse}, so
$a_i=2$.
The branch formula \eqref{eq:branch-count} gives $A_1\geq2-b$.
The inequalities $m_{\mathrm{tr}}\geq M_{\boldsymbol n}=2A_2+A_1$
and $2(m_{\mathrm{tr}}-1)\leq6A_2+A_1-b$ give $A_1\leq k$.
Substituting $2A_2=k+b-2$ yields
\begin{equation}
 2-b\leq A_1\leq k,\qquad
 k\leq m_{\mathrm{tr}}
 \leq\frac{3k+2b-4+A_1}{2}\leq2k.
 \label{eq:d2-preliminary-bounds}
\end{equation}
Every $F\in B$ has a unique decomposition
\[
 F=A_F(v,y)+xB_F(v,y),\qquad
 B_F(v,y)=\sum_{j\geq0}b_j(v)y^j,
\]
with $A_F,B_F\in\C\{v,y\}$.  The summands $A_F(v,y)$ and
$xB_F(v,y)$ are respectively its \emph{even part} and
\emph{odd part} with respect to $x\mapsto-x$.
Indeed, in the power series expansion of $F$, group the even powers of $x$
and factor one $x$ from the odd powers.  The two remaining series lie in
$\C\{x^2,y\}=\C\{v,y\}$, since $x^2=v-y^k$.
Uniqueness follows because the even and odd powers of $x$
cannot cancel.  Thus
$B=\C\{v,y\}\oplus x\C\{v,y\}$ as a module over $\C\{v,y\}$.
The points $(x,y)$ and $(-x,y)$ have the same values of $v,y$.
The even part takes the same value at these points, and the difference
of their $F$-values is $2xB_F(v,y)$.  For pullbacks of target coordinates,
the Jacobian divisibility in \eqref{eq:input-divisor}
constrains the odd part of $F$, and \eqref{eq:d2-preliminary-bounds}
leaves at most two nonzero coefficients of $B_F$ below order $m_{\mathrm{tr}}$ in $y$.
In derivatives such as $(A_F)_y$ and $(B_F)_y$, the variable $v$
is held fixed.

\begin{lemma}
\label{lem:d2-recurrence}
Let $F$ be the pullback under $\boldsymbol n$ of a target coordinate
after a biholomorphic change of target coordinates preserving $v$.
There are $\alpha_F,\beta_F\in\C\{v\}$, with $\alpha_F(0)=0$, such that
\begin{equation}
 F\equiv
 \begin{cases}
  \alpha_F(v)+x\beta_F(v),&m_{\mathrm{tr}}=k,\\
  \alpha_F(v)+x(2v+y^k)\beta_F(v),&k<m_{\mathrm{tr}}\leq2k,
 \end{cases}
 \pmod{y^{m_{\mathrm{tr}}}}.
 \label{eq:d2-A}
\end{equation}
In the first case, $\beta_F(0)=0$.
\end{lemma}

\begin{proof}
Changes of target coordinates preserve \eqref{eq:input-divisor},
since the transformed minors are analytic linear combinations
of the original minors.  Hence
\[
 J(v,F)=2x(A_F)_y+2(v-y^k)(B_F)_y-ky^{k-1}B_F
 \in(y^{m_{\mathrm{tr}}-1}).
\]
Separating even and odd powers of $x$ gives
$(A_F)_y\in(y^{m_{\mathrm{tr}}-1})$ and
\[
 2vjb_j=(2j-k)b_{j-k}\quad(1\leq j<m_{\mathrm{tr}}),
 \qquad b_j=0\quad(j<0).
\]
By \eqref{eq:d2-preliminary-bounds}, $k\leq m_{\mathrm{tr}}\leq2k$,
so all coefficients $b_j$ with
$0\leq j<m_{\mathrm{tr}}$ vanish except possibly $b_0$ and, when
$m_{\mathrm{tr}}>k$, $b_k$.  In the latter case the recurrence at $j=k$
is $b_0=2vb_k$, giving the second congruence.
Integrating $(A_F)_y\in(y^{m_{\mathrm{tr}}-1})$ with respect to $y$ gives
$A_F\equiv\alpha_F(v)\pmod{y^{m_{\mathrm{tr}}}}$.
Finally, $F(0)=0$ gives $\alpha_F(0)=0$, and in the first case
$(dF)_0=0$ also gives $\beta_F(0)=0$.
\end{proof}

\Cref{lem:d2-recurrence} bounds the orders of pullbacks of target
coordinates along the branches of $V$.  For a branch $C$ of
$V$ with normalization parametrization $\gamma_C:(\C,0)\to(C,0)$,
write $\nu_C(F)=\ord_t(F\circ\gamma_C)$.
We set $\nu_C(F)=\infty$ when $F$ vanishes identically on $C$.

\begin{corollary}
\label{cor:d2-branch-orders}
For every $F$ as in \Cref{lem:d2-recurrence} and every branch $C$ of $V$,
\begin{equation}
 b\nu_C(F)\geq\min\{2m_{\mathrm{tr}},3k\}.
 \label{eq:d2-branch-orders}
\end{equation}
\end{corollary}

\begin{proof}
On the normalization of each branch $C$ of $V$,
$\nu_C(y)=2/b$ and $\nu_C(x)=k/b$.
If $m_{\mathrm{tr}}=k$, the first congruence in \eqref{eq:d2-A}
shows that $F$ restricts to zero modulo $y^k$ on $V$, so $b\nu_C(F)\geq2k$.
Otherwise, the second congruence shows that $F$ restricts to
$\beta_F(0)xy^k$ modulo $y^{m_{\mathrm{tr}}}$ on $V$.
The monomial $xy^k$ has order $3k/b$, and the remainder has order at least
$2m_{\mathrm{tr}}/b$.
\end{proof}

We compare the orders along the branches in \eqref{eq:d2-branch-orders} with
$i_0(v,u)=m_X$, from \eqref{eq:input-point}.

\begin{proposition}
\label{prop:d2-numerical}
The integer $k$ is even, and
\[
 (A_1,m_{\mathrm{tr}})=(0,k)\qquad\text{or}\qquad
 (A_1,m_{\mathrm{tr}})=(k,2k).
\]
\end{proposition}

\begin{proof}
Summing \eqref{eq:d2-branch-orders} for $F=u$ over the branches of $V$ gives
\[
 m_X\geq\min\{2m_{\mathrm{tr}},3k\}.
\]
On the other hand, \eqref{eq:d2-counts} and
\eqref{eq:d2-preliminary-bounds} give
\[
 2m_{\mathrm{tr}}\geq2M_{\boldsymbol n}=m_X+A_1,\qquad
 m_X=2k+2b-4+A_1\leq3k.
\]
If $A_1>0$, then $2m_{\mathrm{tr}}>m_X$, so the lower bound
forces $m_X=3k$.  Since $b\leq2$ and $A_1\leq k$, the equality $m_X=3k$ forces
$b=2$ and $A_1=k$.  Now
$2k=M_{\boldsymbol n}\leq m_{\mathrm{tr}}\leq2k$.

If $A_1=0$, the bound $A_1\geq2-b$ gives $b=2$, hence $m_X=2k$.
The lower bound for $m_X$, together with $m_{\mathrm{tr}}\geq k$,
then forces $m_{\mathrm{tr}}=k$.
In both cases $b=2$, so $k=2A_2$ is even.
\end{proof}

\subsection{Exclusion of the second alternative}

For $(A_1,m_{\mathrm{tr}})=(k,2k)$,
\Cref{lem:d2-recurrence} determines the initial forms of the first two
coordinate functions.  These forms are obtained from $(\Phi,\Psi)$ in
\eqref{eq:d2-cubic-pair} by setting $s=x$ and $\lambda=iy^{k/2}$.
We first compare the initial form of the third coordinate at $(x,y)$ and
$(x,\zeta y)$, where $\zeta^k=1$, and then use
\Cref{lem:d2-cubic-orbit}.

\begin{proposition}
\label{prop:d2-second-alternative}
The alternative $(A_1,m_{\mathrm{tr}})=(k,2k)$ is impossible.
\end{proposition}

\begin{proof}
Write $k=2a$ and give $x,y$ weights $a,1$.
The recurrence in \Cref{lem:d2-recurrence} gives, for the pullback $F$ of every target coordinate,
\[
 F\equiv\alpha_F(v)+\beta_F(v)\Theta\pmod{y^{2k}},
 \qquad \Theta=x(2v+y^k).
\]
The coefficient $\beta_u(0)$ is nonzero, since otherwise $u$ would have order
at least $2k$ on each of the two branches of $V$.
This would give $i_0(v,u)\geq4k$, contrary to $i_0(v,u)=m_X=3k$.
Set $U=(u-\alpha_u(v))/\beta_u(v)$ and replace the third coordinate
$w$ by $Z=w-\alpha_w(v)-\beta_w(v)U$.  The map then has the form
\[
 \boldsymbol n=(G,U,Z)
 =\bigl(v,\Theta+y^{2k}C_2,y^{2k}C_3\bigr),
 \qquad C_2,C_3\in B.
\]
The initial forms of $G$ and $U$ for weights $(a,1)$ are, respectively,
\[
 v_0=x^2+y^k,\qquad \Theta_0=2x^3+3xy^k.
\]
They have only the origin as a common zero, and $i_0(G,U)=m_X=3k>1$.
By injectivity and \Cref{lem:filtered-escape}, subtracting a polynomial
in $G,U$ from the third coordinate gives a third coordinate whose initial form
$R\notin\C[v_0,\Theta_0]$ has weight $\rho\geq2k=4a$.

Write $\mu_k=\{\zeta\in\C:\zeta^k=1\}$ for the group of $k$th
roots of unity.  Use $(x,y)$ and $(X,Y)$ for the two source points in the
fiber product of $(v_0,\Theta_0)$.  For $\zeta\in\mu_k$, the graph $X=x,Y=\zeta y$
lies in this fiber product.  Since
$J(v_0,\Theta_0)=-3ky^{2k-1}$, the graph is the unique local component
of this fiber product at each of its points with $y\neq0$.
Consequently, if the polynomial
\[
 R(x,\zeta)-R(x,1)
\]
is nonzero, each of its zeros gives an isolated collision
of $(v_0,\Theta_0,R)$ in the chart $y=1$.
By \Cref{lem:weighted-lifting}, injectivity forces this polynomial to
be constant for every $\zeta$.  Comparing the coefficients of
positive powers of $x$ shows that the exponent of $y$ in each monomial
of $R$ involving $x$ is divisible by $k$.  Hence
\begin{equation}
 R=U_0(x,y^k)+cy^\rho,                              \label{eq:d2-fourier}
\end{equation}
where $U_0\in\C[x,t]$ is homogeneous for weights $a,k$ on $x,t$, respectively.

If $a\mid\rho$, substituting $s=x,\lambda=iy^a$ gives
$v_0=\Phi(s,\lambda)$, $\Theta_0=\Psi(s,\lambda)$, and
$R=\widehat R(s,\lambda)$, where $\widehat R$ is homogeneous
of degree $\rho/a\geq4$.
Since $\widehat R\notin\C[\Phi,\Psi]$,
\Cref{lem:d2-cubic-orbit} gives an isolated collision
of $(\Phi,\Psi,\widehat R)$.
The map
\[
 \mathbb P(a,1,a,1)\longrightarrow\mathbb P^3,\qquad
 [x:y:X:Y]\longmapsto[x:iy^a:X:iY^a]
\]
is well defined since its four coordinate polynomials have weighted degree $a$.
It is finite because the corresponding extension of homogeneous coordinate rings
adjoins $a$-th roots of the second and fourth target coordinates divided
by $i$, and surjective because these roots exist over $\C$.
The difference equations for
$(v_0,\Theta_0,R)$ are pullbacks of those for $(\Phi,\Psi,\widehat R)$,
so $\mathcal Z(v_0,\Theta_0,R)$ is the inverse image of
$\mathcal Z(\Phi,\Psi,\widehat R)$.
The inverse image of the isolated collision is nonempty and consists
of isolated points outside the diagonal, so
\Cref{lem:weighted-lifting} contradicts injectivity.

If $a\nmid\rho$, homogeneity in \eqref{eq:d2-fourier} forces
$R=cy^\rho$, with $c\neq0$.  Choose $\xi^2=-3/2$.
The points $(\xi,1)$ and $(-\xi,1)$ have the same image under
$(v_0,\Theta_0,R)$.  In the chart $y=1$, take the differences
between the values of $v_0,\Theta_0,R$ at $(x,1)$ and $(X,Y)$.
The Jacobian matrix of these differences with respect to $x,X,Y$,
evaluated at $(x,X,Y)=(\xi,-\xi,1)$, is
\[
 \begin{pmatrix}
  2\xi&2\xi&-k\\
  -6&6&3k\xi\\
  0&0&-c\rho
 \end{pmatrix},
\]
with determinant $-24c\rho\xi\neq0$.  The collision is isolated,
and \Cref{lem:weighted-lifting} again contradicts injectivity.
\end{proof}

\subsection{The form of the map along
\texorpdfstring{$y=0$}{y=0}}

In the remaining alternative $(A_1,m_{\mathrm{tr}})=(0,k)$,
the congruence \eqref{eq:d2-A} gives
$F(x,y)\equiv F(x,0)\pmod{y^k}$ for every coordinate function $F$ of $\boldsymbol n$.
We change target coordinates to simplify the restriction of
$\boldsymbol n$ to $y=0$ and obtain the form used in the final collision argument.

\begin{proposition}
\label{prop:d2-remaining-form}
There are an odd integer $p\geq3$ and source and target coordinates
in which
\begin{equation}
 \boldsymbol n(x,y)=
 \bigl(x^2+y^k,x^p+y^kh_2(x,y),y^kh_3(x,y)\bigr),
 \label{eq:d2-remaining-form}
\end{equation}
where $k\geq4$ is even and $h_2,h_3\in\C\{x,y\}$ satisfy
\[
 (h_2(0),h_3(0))\neq(0,0).
\]
\end{proposition}

\begin{proof}
\Cref{prop:d2-numerical,prop:d2-second-alternative} give
$m_{\mathrm{tr}}=k\geq4$ and $m_X=2k$.
\Cref{lem:d2-recurrence} gives
$F(x,y)\equiv\alpha_F(v)+x\beta_F(v)\pmod{y^k}$.  Since
$v\equiv x^2\pmod{y^k}$, it follows that
$F(x,y)\equiv F(x,0)\pmod{y^k}$.

The restriction $\gamma=\boldsymbol n|_{y=0}$ is injective and has
first coordinate $x^2$.  Subtract from the other two coordinates of
$\boldsymbol n$ their even parts on $y=0$, expressed as functions of the first
coordinate $v$.  At least one of the other two coordinate functions
restricts to a nonzero odd function on $y=0$, since $\gamma$ is injective.
Choose as the second coordinate one whose restriction to $y=0$ has
least order $p\geq3$, necessarily odd.  This restriction is $x^p\eta(x^2)$, with
$\eta(0)\neq0$.  Dividing the second coordinate by $\eta$ evaluated
at the first coordinate makes its restriction equal to $x^p$.  By minimality of $p$, the
restriction of the third coordinate is $x^p\xi(x^2)$ for a convergent series $\xi$.
Subtracting $\xi(v)$ times the second coordinate from the third gives
$\gamma=(x^2,x^p,0)$.  The congruences modulo $y^k$
are preserved under these target changes.
The transformed coordinates therefore satisfy \eqref{eq:d2-remaining-form}.

If $h_2(0)=h_3(0)=0$, the last two coordinate functions of $\boldsymbol n$ in
\eqref{eq:d2-remaining-form} would have order greater than $k$ on
each branch of $V$, since $\nu_C(x^p)=kp/2>k$ and
$\nu_C(y^kh_j)>k$.  The original $u$, an analytic function of the
new coordinates vanishing at the origin, would also have order greater than $k$
on each branch of $V$,
because the first coordinate vanishes on $V$.
Then $i_0(v,u)>2k=m_X$, contrary to \eqref{eq:input-point}.
\end{proof}
\subsection{The coordinate \texorpdfstring{$Q=x^p+y^k\beta(y)$}{Q}
and the weighted initial form of \texorpdfstring{$Z$}{Z}}

By \Cref{prop:d2-remaining-form},
\begin{equation}
 \boldsymbol n=(P,Q_0,Z_0)
 =\bigl(x^2+y^k,x^p+y^kh_2,y^kh_3\bigr),\qquad
 (h_2(0),h_3(0))\neq(0,0),                         \label{eq:d2-cyc-start}
\end{equation}
where $k\geq4$ is even, $p\geq3$ is odd, and
$h_2,h_3\in\C\{x,y\}$.  To study the fiber product of the first two
coordinates in Step~\ref{step:third-coordinate}, we seek coordinates
in which
\begin{equation}
 (P,Q,Z)=\bigl(x^2+y^k,x^p+y^k\beta(y),Z(x,y)\bigr),
 \qquad \beta\in\C\{y\},\quad Z\in\C\{x,y\},\quad \beta(0)=c\neq0,       \label{eq:d2-cyc-prepared}
\end{equation}
with
\begin{equation}
 \operatorname{in}_{(k,p)}Z=x^N,
 \qquad N\geq3\text{ odd}.                         \label{eq:d2-cyc-xN}
\end{equation}
In these coordinates, the third coordinate $Z$ must distinguish
the points in each fiber of the finite map $(P,Q)$.  Thus any distinct
points with equal $P,Q,Z$ values contradict injectivity, as
in \eqref{eq:strategy-separation}.

\begin{proposition}\label{prop:d2-cyclic-preparation}
Every injective germ in \eqref{eq:d2-cyc-start} can be transformed
by biholomorphic changes of source and target coordinates into a germ
satisfying \eqref{eq:d2-cyc-prepared} and \eqref{eq:d2-cyc-xN}.
\end{proposition}

\begin{proof}
\noindent\emph{Putting $Q$ in the form $x^p+y^k\beta(y)$.}
For a suitable $\alpha\in\C$, the linear change $Q_0\mapsto Q_0+\alpha Z_0$
gives $Q=x^p+y^kh(x,y)$ with $h(0)=c\neq0$.  For an unknown convergent
series $\theta\in\C\{x,y\}$, set
\[
 A_\theta=1-2x\theta-y^k\theta^2,\qquad
 X_\theta=x+y^k\theta,\qquad Y_\theta=yA_\theta^{1/k},
\]
using the $k$th root of $A_\theta$ with constant term one.  This
substitution preserves $P$ and the ideal $(y^k)$, and fixes the axis
$y=0$ pointwise.  Moreover,
\[
 Q(X_\theta,Y_\theta)=x^p+y^k\mathcal H(\theta),\qquad
 \mathcal H(\theta)=
 \frac{(x+y^k\theta)^p-x^p}{y^k}
 +A_\theta h(X_\theta,Y_\theta).
\]
The quotient in this expression is analytic.  We seek an analytic
function $\theta(x,y)$ for which $\mathcal H(\theta(x,y))$ depends only on $y$.
Regard the argument of $\mathcal H$ temporarily as an independent
variable $t$.  At $y=0$,
\[
 \mathcal H(x,0,t)
 =h(x,0)+t\bigl(px^{p-1}-2xh(x,0)\bigr).
\]
Put $t_0=h_x(0)/(2c)$.  The point $(x,t)=(0,t_0)$ is critical
for $\mathcal H(x,0,t)$, and its Hessian in $(x,t)$ has determinant
$-4c^2\neq0$.  The implicit-function theorem gives an analytic
critical point $(a(y),b(y))$ of $(x,t)\mapsto\mathcal H(x,y,t)$
with $(a(0),b(0))=(0,t_0)$.
Let $\beta(y)=\mathcal H(a(y),y,b(y))$ be the corresponding
critical value, so $\beta(0)=c$.
For each fixed $y$ near zero, the level curve
$\mathcal H(x,y,t)=\beta(y)$ has two transverse analytic branches
at $(a(y),b(y))$.
We choose a branch of the form $t=\theta(x,y)$ by the following
factorization.

Write $\xi=x-a(y)$ and $\tau=t-b(y)$.  Taylor expansion at the
critical point gives
\[
 \mathcal H-\beta=a_{20}\xi^2+2a_{11}\xi\tau+a_{02}\tau^2,
 \qquad a_{11}(0)=-c,\quad a_{02}(0)=0,
\]
where $a_{20},a_{11},a_{02}$ are analytic in $(\xi,\tau,y)$.
Choose the analytic square root $S$ of $a_{11}^2-a_{20}a_{02}$
with $S(0)=-c$.  Since $a_{11}+S$ is a unit,
\[
 \mathcal H-\beta=
 \bigl(a_{20}\xi+(a_{11}+S)\tau\bigr)
 \left(\xi+\frac{a_{02}}{a_{11}+S}\tau\right).
\]
The first factor has nonzero $t$-derivative at $(x,y,t)=(0,0,t_0)$.
The implicit-function theorem therefore gives an analytic
$t=\theta(x,y)$ on which $\mathcal H=\beta(y)$.
The derivative of $(X_\theta,Y_\theta)$ at the origin is the identity.
This change of coordinates preserves the function $P$ and the ideal
$(y^k)$, and fixes $y=0$ pointwise.

After this substitution, the second coordinate is $x^p+y^k\beta(y)$,
which we again denote by $Q$.
The transformed third coordinate is $\widetilde Z_0=y^k\widetilde h_3$,
with $\widetilde h_3\in\C\{x,y\}$.
Writing $p=2m+1$, we have
\begin{equation}
 Q\equiv y^k\bigl(\beta(y)+(-1)^mxy^{k(m-1)}\bigr)\pmod P,
 \qquad (P,Q)\C\{x,y\}=(x^2,y^k).
 \label{eq:d2-cyc-plane-ideal}
\end{equation}
The factor $\beta(y)+(-1)^mxy^{k(m-1)}$ is a unit
because its constant term is $c\neq0$.
The quotient ring $\C\{x,y\}/(x^2,y^k)$ has $\C$-basis
$1,y,\ldots,y^{k-1},x,xy,\ldots,xy^{k-1}$.  Thus $(P,Q)$ is finite of degree
$2k$, by \Cref{lem:finite-plane-degree}, and injectivity implies
$\widetilde Z_0\neq0$.  If $d=\widetilde h_3(0)\neq0$, replacing the
third coordinate by $Q-(c/d)\widetilde Z_0$ gives initial form
$x^p$ for weights $(k,p)$, as required.
For the remaining case, write the transformed third coordinate as
$Z_0=y^kh$, with $0\neq h\in(x,y)$.

\smallskip
\noindent\emph{Subtracting functions of $P,Q$.}
Use weights $(k,p)$ on $(x,y)$, and put
\[
 P_{\mathrm{in}}=x^2,\quad Q_{\mathrm{in}}=x^p+cy^k,
 \quad A=\C\{P,Q\},\quad \mathcal R=Q^2-P^p.
\]
The polynomial map $(P_{\mathrm{in}},Q_{\mathrm{in}})$ is finite, so
\Cref{lem:filtered-escape} applies and
$\operatorname{gr}A=\C[P_{\mathrm{in}},Q_{\mathrm{in}}]$.
Identify $\C[S,T]$ with $\C[P_{\mathrm{in}},Q_{\mathrm{in}}]$ by
$S\mapsto P_{\mathrm{in}}$, $T\mapsto Q_{\mathrm{in}}$.
Restriction to $y=0$ is the homomorphism
\[
 \C[S,T]\longrightarrow\C[x],\qquad
 S\longmapsto x^2,\quad T\longmapsto x^p.
\]
Division of a polynomial in $\C[S,T]$ by $T^2-S^p$, monic in $T$,
leaves a remainder $A_0(S)+TA_1(S)$.  Its image is
$A_0(x^2)+x^pA_1(x^2)$.
Since $p$ is odd, the first summand has only even powers of $x$
and the second only odd powers.  Their sum vanishes only when
$A_0=A_1=0$.  The kernel of the restriction homomorphism is therefore
$(Q_{\mathrm{in}}^2-P_{\mathrm{in}}^p)$ in $\C[P_{\mathrm{in}},Q_{\mathrm{in}}]$.
Expansion of $Q^2-P^p$ gives $\mathcal R\in y^k(x,y)$ and
$\operatorname{in}_{(k,p)}\mathcal R=Q_{\mathrm{in}}^2-P_{\mathrm{in}}^p$.
Starting with $Z_0$, whenever the current remainder has initial form
in $\operatorname{gr}A$, we cancel that initial form by subtracting
$\mathcal R C(P,Q)$ for a polynomial $C(S,T)$ weighted homogeneous
with $\deg S=2k$ and $\deg T=pk$.
These subtractions preserve $y^k(x,y)$.
The proof of \Cref{lem:filtered-escape} shows that they terminate unless
$Z_0\in A$, which injectivity excludes by \eqref{eq:d2-cyc-plane-ideal}.
The remainder obtained from $Z_0$ therefore has initial form
\begin{equation}
 H=y^k\Lambda\notin\C[P_{\mathrm{in}},Q_{\mathrm{in}}],            \label{eq:d2-cyc-firstH}
\end{equation}
where $\Lambda$ is weighted homogeneous of positive degree.

In this proof, a polynomial is invariant if it is unchanged by every
substitution $(x,y)\mapsto(x,\zeta y)$ with $\zeta\in\mu_k$.
A monomial $x^iy^j$ is invariant under this action exactly when $k\mid j$.
For $\zeta\in\mu_k\setminus\{1\}$, let
$D_\zeta(s)=\Lambda(s,1)-\Lambda(s,\zeta)$.
A nonzero root $s_0$ of a nonzero $D_\zeta$ would give an isolated
collision of $(P_{\mathrm{in}},Q_{\mathrm{in}},H)$ at
$(s_0,1),(s_0,\zeta)$, since locally the first two
difference equations select $X=x$, $Y=\zeta y$, and the difference
of the $H$ values restricts to $D_\zeta$ after setting $y=1$.
Consequently, injectivity and \Cref{lem:weighted-lifting} imply that
a nonzero $D_\zeta\in\C[s]$ has no root outside $s=0$.
Factoring this polynomial over $\C$ shows that it is a nonzero
scalar multiple of a power of $s$, allowing exponent zero.
Thus every $D_\zeta$ is either identically zero or a monomial.
When $\zeta$ is a primitive $k$th root of unity, distinct noninvariant monomials of
$\Lambda$ give distinct powers of $s$, by weighted homogeneity.
Hence either $\Lambda\in\C[x,y^k]$, or
\begin{equation}
 \Lambda=\Lambda_{\mathrm{inv}}+\gamma x^iy^j,
 \quad \Lambda_{\mathrm{inv}}\in\C[x,y^k],\quad
 \gamma\neq0,\quad i,j\geq0,\quad k\nmid j.            \label{eq:d2-cyc-mixedLambda}
\end{equation}

We exclude \eqref{eq:d2-cyc-mixedLambda} by considering pairs
$(1,y),(-1,Y)$ with equal values of $P_{\mathrm{in}},Q_{\mathrm{in}}$.
These equalities define the affine curve
\[
 C=\{(y,Y)\in\C^2:Y^k-y^k=2/c\}.
\]
Using $Y^k=y^k+2/c$, write
$y^k\Lambda_{\mathrm{inv}}(1,y)-Y^k\Lambda_{\mathrm{inv}}(-1,Y)$
as $B(y^k)$ on $C$.
Let $\mathcal G\in\C[y,Y]$ be the polynomial on the right below.
On $C$,
\[
 \mathcal G(y,Y)=H(1,y)-H(-1,Y)
 =B(y^k)+\gamma\bigl(y^{k+j}-(-1)^iY^{k+j}\bigr),
 \qquad B\in\C[z].
\]
To obtain the collision required in Step~\ref{step:third-coordinate},
we show that $\mathcal G$ has an isolated zero on $C$.
We prove that its homogenization has a zero on the projective closure
of $C$ away from infinity.

The projective closure of $C$,
\[
 \overline C=\{[y:Y:T]\in\mathbb P^2:Y^k-y^k=(2/c)T^k\},
\]
is smooth because the three partial derivatives of its defining
polynomial have no common projective zero.  It is irreducible because
distinct projective plane components would meet and produce a
singular point.  Let $\mathcal G_{\mathrm{hom}}$ be the
homogenization of $\mathcal G$ with respect to $T$.  The line at
infinity is $T=0$.
If $B\neq0$ and $k\deg B>k+j$, then
$\mathcal G_{\mathrm{hom}}$ equals the nonzero leading coefficient
of $B$ at every point $[1:\zeta:0]$ of $\overline C$ at infinity.
B\'ezout's theorem \cite[Section~2.6, p.~54]{CasasAlvero}
then gives an affine zero of $\mathcal G|_C$,
since $\overline C$ and $V(\mathcal G_{\mathrm{hom}})$ have positive
degrees and no intersection at infinity.

Otherwise, $k\nmid j$ implies that $\mathcal G$ has total degree $k+j$.
At infinity its homogenization has values
$\gamma(1-(-1)^i\zeta^{k+j})$, where $\zeta^k=1$.
These values are not all zero, since $k\nmid j$.
Thus $\mathcal G_{\mathrm{hom}}$ does not vanish identically on
$\overline C$, and the total intersection number of
$V(\mathcal G_{\mathrm{hom}})$ with $\overline C$ is $k(k+j)$.
At a zero of $\mathcal G_{\mathrm{hom}}$ on $\overline C$ at infinity,
use $T$ as local parameter on $\overline C$ in the chart $y=1$:
\[
 Y=\zeta(1+2T^k/c)^{1/k},\qquad
 \zeta^k=1,\qquad (-1)^i\zeta^{k+j}=1.
\]
The contribution of $\gamma\bigl(y^{k+j}-(-1)^iY^{k+j}\bigr)$
to $\mathcal G_{\mathrm{hom}}$ in this parametrization is
\[
 \gamma\bigl(1-(1+2T^k/c)^{(k+j)/k}\bigr),
\]
which has order exactly $k$.  The contribution of $B(y^k)$ has
$T$-exponents congruent to $j$ modulo $k$, so it cannot cancel
the coefficient of $T^k$.  The intersection multiplicity at each of
the $k$ points at infinity is therefore at most $k$.  Since $j>0$,
the sum of these multiplicities satisfies
\[
 \sum_{a\in\overline C\cap\{T=0\}}
 i_a(\overline C,V(\mathcal G_{\mathrm{hom}}))
 \leq k^2<k(k+j).
\]
Hence an affine zero of $\mathcal G|_C$ exists in this case as well.

In both cases, the restriction of $\mathcal G$ to the irreducible
curve $C$ does not vanish identically, so its zeros are isolated.  Any such zero
gives distinct points $(1,y),(-1,Y)$ with equal values of
$P_{\mathrm{in}},Q_{\mathrm{in}},H$.
Setting $x=1$ gives a finite cover of the weighted
projective chart $x\neq0$.  On this cover, the first difference
equation locally selects $X=-1$, and the remaining equations
are those of $C\cap V(\mathcal G)$.  The collision is therefore
isolated in $\mathcal Z(P_{\mathrm{in}},Q_{\mathrm{in}},H)$.
By \Cref{lem:weighted-lifting}, this contradicts injectivity.
Therefore $\Lambda\in\C[x,y^k]$.

\smallskip
\noindent\emph{The case $\Lambda\in\C[x,y^k]$.}
Set $t=y^k$ and give $x,t$ degrees $1,p$.  Since
$t=c^{-1}(Q_{\mathrm{in}}-xP_{\mathrm{in}}^{(p-1)/2})$,
\[
 \C[x,t]=\C[P_{\mathrm{in}},Q_{\mathrm{in}}]
       \oplus x\C[P_{\mathrm{in}},Q_{\mathrm{in}}].
\]
The sum is direct because the involution
$(x,t)\mapsto(-x,t+2x^p/c)$ fixes $P_{\mathrm{in}},Q_{\mathrm{in}}$
and sends $x$ to $-x$.
Write
\[
 H=A_{\mathrm{ev}}(P_{\mathrm{in}},Q_{\mathrm{in}})
       +xA_{\mathrm{odd}}(P_{\mathrm{in}},Q_{\mathrm{in}}).
\]
By \eqref{eq:d2-cyc-firstH}, $A_{\mathrm{odd}}\neq0$.
If $A_{\mathrm{odd}}$ is not a nonzero multiple of a power of
$P_{\mathrm{in}}$, weighted homogeneity makes
$A_{\mathrm{odd}}(1,q)$ a nonconstant polynomial in $q$.
Choose a root $q_0$ and choose $y,Y$ satisfying
\[
 y^k=(q_0-1)/c,\qquad Y^k=(q_0+1)/c.
\]
The distinct points $(1,y),(-1,Y)$ have equal values of
$P_{\mathrm{in}},Q_{\mathrm{in}},H$.
The collision is isolated.  After setting $x=1$, the first
difference equation fixes $X=-1$, the second gives
$Y^k-y^k=2/c$, and the difference of the $H$ values restricts to
the nonzero polynomial
$2A_{\mathrm{odd}}(1,1+cy^k)$.
There are only finitely many solutions for $y$ and $Y$, including
when one of them is zero.  This contradicts
\Cref{lem:weighted-lifting}.
Therefore $A_{\mathrm{odd}}=\kappa_0P_{\mathrm{in}}^{r_0}$
with $\kappa_0\neq0$.
Subtracting $A_{\mathrm{ev}}(P,Q)$ from the remainder with initial form $H$
and dividing by $\kappa_0$ gives a third coordinate $Z$ with initial form $x^N$
for weights $(k,p)$, where $N=2r_0+1$ is odd.
Since $H=y^k\Lambda$ and $\Lambda$ has positive weight for
$\deg x=1$, $\deg(y^k)=p$, one has $N>p\geq3$.
\end{proof}

\subsection{A collision on a component of the fiber product}

Assume that $\boldsymbol n=(P,Q,Z)$ satisfies
\eqref{eq:d2-cyc-prepared} and \eqref{eq:d2-cyc-xN}.
We use the weighted initial form of $Z$ to study the difference
$Z(x,y)-Z(X,Y)$ on a local component of the fiber product of $(P,Q)$.
As in Step~\ref{step:third-coordinate}, injectivity requires this
difference to be nonzero at every pair of distinct points.
If $\beta(\zeta y)=\beta(y)$ for every $\zeta\in\mu_k$,
Weierstrass preparation forces a collision.  Otherwise the term of
least degree in $\beta$ whose exponent is not divisible by $k$
determines a second choice of weights,
for which injectivity implies a polynomial identity
that forces a collision.

\begin{proposition}\label{prop:d2-cyclic-separation}
No germ satisfying \eqref{eq:d2-cyc-prepared} and
\eqref{eq:d2-cyc-xN} is injective.
\end{proposition}

\begin{proof}
Suppose otherwise, and fix a primitive $k$th root of unity $\zeta$.
Write $\beta(y)=\sum_{h\geq0}\beta_hy^h$, with $\beta_0=c$.

\smallskip
\noindent\emph{The local fiber product.}
Put $\nu=k(p-2)$ and consider
\[
 P_\lambda=x^2+\lambda^\nu y^k,\qquad
 Q_\lambda=x^p+y^k\beta(\lambda^py),\qquad
 R_\lambda=\lambda^{-kN}Z(\lambda^kx,\lambda^py).
\]
These functions are analytic, with $R_0=x^N$ by
\eqref{eq:d2-cyc-xN}, and for
$\lambda\neq0$ the map $(P_\lambda,Q_\lambda,R_\lambda)$ is obtained
from the original germ by invertible scalings.  Set $y=1$ and compare $(s,1),(S,z)$
near $(0,1),(0,\zeta)$.  The equations for the fiber product of
$(P_\lambda,Q_\lambda)$ are
\begin{equation}
 \begin{aligned}
 s^2-S^2+\lambda^\nu(1-z^k)&=0,\\
 s^p-S^p+\beta(\lambda^p)-z^k\beta(\lambda^pz)&=0.
 \end{aligned}
 \label{eq:d2-cyc-fiber-product}
\end{equation}
The derivative with respect to $z$ of the left-hand side of the second
equation at $(s,S,z,\lambda)=(0,0,\zeta,0)$ is
$-ck\zeta^{k-1}\neq0$.  Thus the second equation has a unique analytic
solution
$z=\mathfrak z(\upsilon,\lambda)$, where $\upsilon=s^p-S^p$.
Write
\[
 1-\mathfrak z(\upsilon,\lambda)^k
 =d_0(\lambda)+\upsilon\mathcal K(\upsilon,\lambda),\qquad
 K_p=(s^p-S^p)/(s-S).
\]
The quotient $K_p$ is a polynomial in $s,S$.
The first equation becomes
\begin{equation}
 (s-S)\bigl(s+S+\lambda^\nu K_p\mathcal K\bigr)
            +\lambda^\nu d_0(\lambda)=0.             \label{eq:d2-cyc-hypersurface}
\end{equation}
The functions $U=s-S$,
$\Xi=s+S+\lambda^\nu K_p\mathcal K$, and $\lambda$ are analytic
coordinates near $(s,S,\lambda)=(0,0,0)$, since their Jacobian
determinant there is nonzero.
On this local fiber product, let $D=R_\lambda(s,1)-R_\lambda(S,z)$.
Since $z$ remains near $\zeta\neq1$, this neighborhood of the
fiber product is disjoint from the diagonal.
In particular, $U=0$ does not describe coincident source points.
Injectivity requires the zero set of $D$ to be contained in
the set $\lambda=0$.

If $\beta(\zeta y)=\beta(y)$, uniqueness gives
$\mathfrak z(0,\lambda)=\zeta$, so $d_0=0$.
On the component $\Xi=0$ of the fiber product $U\Xi=0$,
the restriction of $D$ at $\lambda=0$ is
$D|_{\lambda=0}=2^{1-N}U^N$.  Weierstrass preparation
\cite[Theorem~I.1.6]{GreuelLossenShustin} gives zeros of $D$ near
$U=0$ for every small $\lambda\neq0$.  Since $z$ remains near
$\zeta\neq1$, these are collisions of distinct points.
We may therefore assume $\beta(\zeta y)\neq\beta(y)$.

Let $j$ be the least index with $\beta_j\neq0$ and $\zeta^j\neq1$,
and put
\[
 L=k+j,\qquad \chi=\zeta^j,\qquad
 T=-\frac{\beta_j}{c}y^L,\qquad \rho=-c(x^2+T).
\]
Comparing coefficients in
$\beta(q)=\mathfrak z(0,\lambda)^k\beta(q\mathfrak z(0,\lambda))$,
where $q=\lambda^p$, gives
$\mathfrak z(0,\lambda)=\zeta+O(q^j)$, since all terms of degree
less than $j$ in $\beta$ are invariant under $y\mapsto\zeta y$.
The coefficient of $q^j$ then gives
\begin{equation}
 d_0(\lambda)=-\frac{\beta_j}{c}(1-\chi)\lambda^{pj}
            +O(\lambda^{p(j+1)}).                    \label{eq:d2-cyc-smoothing}
\end{equation}
Set $r=\nu+pj=pL-2k$, and write
$\lambda^\nu d_0=\lambda^r\varepsilon(\lambda)$, where
$\varepsilon(0)=-(\beta_j/c)(1-\chi)\neq0$.
Replacing $\Xi$ by $-\Xi/\varepsilon(\lambda)$ puts the local fiber
product in the form $U\Xi=\lambda^r$.
Injectivity excludes zeros of $D$ over $\lambda\neq0$, while
on $\lambda=0$ its restrictions are
$D|_{U=0}=0$, $D|_{\Xi=0}=2^{1-N}U^N$.

The two branches $U=0$ and $\Xi=0$ of the fiber $\lambda=0$
correspond respectively to $s=S$ and $s=-S$ when $\lambda=0$.
On the first branch, $z=\mathfrak z(0,0)=\zeta$.  On the second,
$z=\mathfrak z(2s^p,0)$ satisfies $z^k=1+2s^p/c$ and tends to
$\zeta$ as $s\to0$.
The restrictions of $D$ to the two branches of $\lambda=0$,
together with the absence of zeros when $\lambda\neq0$, determine
its divisor of zeros on $\{U\Xi=\lambda^r\}$.  To see this, pull back by
the finite cover
\[
 \pi:(\C^2,0)\longrightarrow(\{U\Xi=\lambda^r\},0),\qquad
 (\kappa_1,\kappa_2)\longmapsto
 (\kappa_1^r,\kappa_2^r,\kappa_1\kappa_2).
\]
The zeros of $\pi^*D$ lie on the coordinate axes.
Unique factorization in $\C\{\kappa_1,\kappa_2\}$
\cite[Theorem~I.1.16]{GreuelLossenShustin}
and the restriction
\[
 (\pi^*D)|_{\kappa_2=0}=2^{1-N}\kappa_1^{rN}
\]
therefore give $\pi^*D=\kappa_1^{rN}\widetilde\varepsilon_D$
with $\widetilde\varepsilon_D$ a unit.
The quotient $\widetilde\varepsilon_D=\pi^*D/\kappa_1^{rN}$
is invariant under
$(\kappa_1,\kappa_2)\mapsto(\xi\kappa_1,\xi^{-1}\kappa_2)$,
$\xi^r=1$.  A monomial $\kappa_1^a\kappa_2^b$ is invariant exactly
when $r$ divides $a-b$.  Grouping these monomials gives
\[
 \C\{\kappa_1,\kappa_2\}^{\mu_r}
 =\C\{\kappa_1^r,\kappa_2^r,\kappa_1\kappa_2\}
 \cong\C\{U,\Xi,\lambda\}/(U\Xi-\lambda^r).
\]
Thus $\widetilde\varepsilon_D$ is the pullback $\pi^*\varepsilon_D$
of a unit $\varepsilon_D$ on $\{U\Xi=\lambda^r\}$, and
\begin{equation}
 D=U^N\varepsilon_D,
 \qquad \varepsilon_D\in
 \bigl(\C\{U,\Xi,\lambda\}/(U\Xi-\lambda^r)\bigr)^*.             \label{eq:d2-cyc-hyp-descent}
\end{equation}

\smallskip
\noindent\emph{Rescaling with weights $(L,2)$.}
The substitutions $\lambda=\tau^2$, $s=\tau^r\widehat s$, and
$S=\tau^r\widehat S$ give both $s^2-S^2$ and
$\lambda^r\varepsilon(\lambda)$ in \eqref{eq:d2-cyc-hypersurface}
a factor $\tau^{2r}$.
Since $2k+r=pL$, the points in the original source coordinates become
\[
 (\tau^{pL}\widehat s,\tau^{2p}),\qquad
 (\tau^{pL}\widehat S,\tau^{2p}z),
\]
where $(x,y)$ have weights $(L,2)$ and the scaling parameter is $\tau^p$.
After division by $\tau^{2r}$, the term containing
$s^p-S^p$ in \eqref{eq:d2-cyc-hypersurface} has the factor $\tau^{2\nu+rp-2r}=\tau^{p(p-2)L}$.
Since $p\geq3$ and $L>0$, the exponent $p(p-2)L$ is positive,
so this term vanishes at $\tau=0$.
Thus \eqref{eq:d2-cyc-smoothing} gives the curve obtained at $\tau=0$
\begin{equation}
 \widehat S^2-\widehat s^2=(1-\chi)T_0,
 \qquad T_0=-\beta_j/c\neq0.
 \label{eq:d2-cyc-limit-curve}
\end{equation}
At each point $(u,v)$ of this curve,
$v^2-u^2=(1-\chi)T_0\neq0$.
Thus $u,v$ cannot both vanish, and $u\neq v$.
After the substitutions above, the difference of the $Q_\lambda$
values has nonzero derivative in $z$ at
$(\widehat s,\widehat S,z,\tau)=(u,v,\zeta,0)$, while the polynomial
$\widehat S^2-\widehat s^2-(1-\chi)T_0$ defining the limiting curve
\eqref{eq:d2-cyc-limit-curve} has gradient $(-2u,2v)$ in
$(\widehat s,\widehat S)$.
Fix $\widehat s=u$ if $v\neq0$, or $\widehat S=v$ otherwise.
The implicit-function theorem gives an analytic lift of $(u,v)$ to
the rescaled fiber product, represented by functions
$(\widehat s(\tau),\widehat S(\tau),z(\tau))$ satisfying the rescaled
form of \eqref{eq:d2-cyc-fiber-product}, with initial value $(u,v,\zeta)$.
Equation \eqref{eq:d2-cyc-hyp-descent} implies along every lift
\begin{equation}
 \begin{aligned}
 &Z(\tau^{pL}\widehat s,\tau^{2p})
   -Z(\tau^{pL}\widehat S,\tau^{2p}z)\\
 &\qquad=\tau^{pLN}(\widehat s-\widehat S)^N(c_D+O(\tau)),
 \qquad c_D=\varepsilon_D(0)\neq0.                  \label{eq:d2-cyc-primitiveD}
 \end{aligned}
\end{equation}
The coefficient of $\tau^{pLN}$ in this difference is
$c_D(u-v)^N\neq0$, so its order is exactly $pLN$.
The constant $c_D=\varepsilon_D(0)$ is independent of the lift.

Let $C_\zeta\subset\mathbb P(L,2,L,2)$ be the curve defined by
$Y=\zeta y$, $\rho(x,y)=\rho(X,Y)$.
It lies in $\mathcal Z(y^k,\rho)$, and its homogeneous coordinate
ring, after eliminating $Y$, is
\[
 R=\C[x,y,X]/(X^2-x^2-(1-\chi)T).
\]
The ring $R$ is a domain.  Indeed, put $\delta_c=X-x$ and
$\sigma_c=X+x$.  The polynomial defining $R$ becomes
$\delta_c\sigma_c-(1-\chi)T(y)$ in $\C[\delta_c,y][\sigma_c]$.
Its coefficients $\delta_c$ and $-(1-\chi)T(y)$ are coprime in
$\C[\delta_c,y]$, since $T(y)$ is a nonzero scalar multiple of
$y^L$.  Since the polynomial has degree one in $\sigma_c$ and its
coefficients are coprime, Gauss's lemma implies that it is irreducible in
$\C[\delta_c,y,\sigma_c]$.
Thus the chart $y\neq0$ of $C_\zeta$ is dense.  Setting $y=1$ gives
a finite cover of this chart by the curve \eqref{eq:d2-cyc-limit-curve}.

Put $t=y^k$.  The terms of $\beta(y)$ whose exponents are divisible
by $k$ sum to $\overline\beta(y^k)$, where
$\overline\beta\in\C\{t\}$ is characterized by
$\overline\beta(y^k)=k^{-1}\sum_{\eta\in\mu_k}\beta(\eta y)$.
For $Q^\flat=Q-P\overline\beta(P)$, minimality of $j$ gives
\begin{equation}
 \operatorname{in}_{(L,2)}P=t,\qquad
 \operatorname{in}_{(L,2)}Q^\flat=\rho.
 \label{eq:d2-cyc-second-initial}
\end{equation}
Indeed, $L=k+j>k$, so $y^k$ has weight $2k<2L$, the weight of $x^2$.
Since $p\geq3$, the term $x^p$ has weight $pL>2L$.
The term of least degree in $y^k\beta(y)$ whose exponent is not
divisible by $k$ is $\beta_jy^L$,
of weight $2L$, and
$P\overline\beta(P)-t\overline\beta(t)$ has initial form $cx^2$.
The map $(t,\rho)$ is finite, and
$\C\{P,Q^\flat\}=\C\{P,Q\}$.  We may therefore apply
\Cref{lem:filtered-escape} with $\C[t,\rho]$ as the algebra generated
by the initial forms of $P,Q^\flat$.

For this application, let $H\in\C[x,y]$ and suppose
$H(x,y)=H(X,\zeta y)$ in $R$.
The automorphism of $R$ fixing $x,y$ and sending $X$ to $-X$ gives
$H(X,\zeta y)=H(-X,\zeta y)$.
The inclusion $\C[X,y]\to R$ is injective because the
relation defining $R$ can be written as a monic equation in $x$, so
$H$ is even in its first variable.  Write $H(x,y)=G(x^2+T,y)$
with $G\in\C[\mathcal S,y]$.
Since $X^2+\chi T=x^2+T$, the identity $H(x,y)=H(X,\zeta y)$ becomes
\[
 G(x^2+T,y)=G(x^2+T,\zeta y).
\]
The homomorphism
\[
 \C[\mathcal S,y]\longrightarrow\C[x,y],\qquad
 \mathcal S\longmapsto x^2+T,\quad y\longmapsto y
\]
is injective.
Since $\zeta$ is a primitive $k$th root of unity,
$G\in\C[\mathcal S,y^k]$, and hence $H\in\C[\rho,t]$.
We have proved
\begin{equation}
 H(x,y)=H(X,\zeta y)\text{ in }R
 \quad\Longrightarrow\quad H\in\C[t,\rho].         \label{eq:d2-cyc-equalizer}
\end{equation}

By \eqref{eq:d2-cyc-plane-ideal}, the map $(P,Q)$ has degree $2k>1$.
Injectivity therefore excludes $Z\in\C\{P,Q\}$.
Thus \Cref{lem:filtered-escape}, applied to $Z$, gives a remainder
$\widetilde Z$ with
$H=\operatorname{in}_{(L,2)}\widetilde Z\notin\C[t,\rho]$.
Write $\delta=\ord_{(L,2)}\widetilde Z$.
By \eqref{eq:d2-cyc-equalizer}, the difference $H(x,y)-H(X,\zeta y)$
is nonzero at some $(u,1),(v,\zeta)$ on the dense chart of $C_\zeta$.
On the analytic lift of this pair the difference of the $\widetilde Z$
values has order $p\delta$ in $\tau$.  Subtracting functions of $P,Q$
preserves the difference of the $Z$ values on the fiber product of $(P,Q)$, so
\eqref{eq:d2-cyc-primitiveD} gives $\delta=NL$.
We compare the coefficients of $\tau^{pLN}$ along the analytic
lift of every point of \eqref{eq:d2-cyc-limit-curve}.
The polynomials $H(x,y)-H(X,\zeta y)$ and $c_D(x-X)^N$ both have
weight $NL$ for weights $(L,2,L,2)$, so equality for $y=1$ extends to
the dense chart $y\neq0$ and hence to all of $C_\zeta$:
\begin{equation}
 H(x,y)-H(X,\zeta y)=c_D(x-X)^N\quad\text{on }C_\zeta.             \label{eq:d2-cyc-Hidentity2}
\end{equation}

\smallskip
\noindent\emph{The final collision.}
Only the odd part of $H$ in $x$ can distinguish $(x,y)$ from $(-x,y)$.
Since $H$ has weight $NL$ and $N$ is odd by
\eqref{eq:d2-cyc-xN}, every odd power $x^i$
in $H$ has $y$-exponent $L(N-i)/2$, a multiple of $L$.
Thus its odd part can be written
\[
 \frac{H(x,y)-H(-x,y)}2=A_N(x,T),
 \qquad \deg x=1,\quad\deg T=2,\quad\deg A_N=N.
\]
The simultaneous sign change $(x,X)\mapsto(-x,-X)$ preserves
$C_\zeta$ and negates the right-hand side of
\eqref{eq:d2-cyc-Hidentity2}.  Subtracting the transformed identity
from the original gives
\begin{equation}
 A_N(x,T)-A_N(X,\chi T)=c_D(x-X)^N
 \quad\text{when }X^2-x^2=(1-\chi)T.
 \label{eq:d2-cyc-Aidentity}
\end{equation}
Setting $X=-x$ and $T=0$ gives
\[
 A_N(x,0)=2^{N-1}c_Dx^N.
\]
The order $N$ of $c_D(x-X)^N$ along $X=x$ also constrains the partial
derivatives of $A_N$ at $T=0$.  Substitute
\[
 X=x+h,\qquad T=\frac{2xh+h^2}{1-\chi}
\]
and differentiate \eqref{eq:d2-cyc-Aidentity} at $h=0$.
Since $N\geq3$, the derivative of its right-hand side is zero, so
\[
 2x\,\partial_TA_N(x,0)-\partial_xA_N(x,0)=0,
 \qquad
 \partial_TA_N(x,0)=2^{N-2}Nc_Dx^{N-2}.
\]
For $T_0=-\beta_j/c\neq0$, the odd polynomial $A_N(x,T_0)$
therefore has nonzero terms of degrees $N$ and $N-2$.  Dividing by its
lowest power of $x$ leaves a polynomial of positive degree with nonzero
constant term, hence a nonzero root $x_0$.
The initial forms in \eqref{eq:d2-cyc-second-initial} take the same values
at $(x_0,1),(-x_0,1)$, and the difference of the $H$ values
is $2A_N(x_0,T_0)=0$.
After setting $y=1$, the collision is isolated, since the first difference
equation for $(t,\rho,H)$ fixes $Y=1$, the second fixes $X=-x$ near $x_0\neq0$,
and the difference of the $H$ values restricts to the nonzero polynomial
$2A_N(x,T_0)$.
This contradicts \Cref{lem:weighted-lifting}.
\end{proof}

\begin{theorem}
\label{thm:dtwo}
Under the hypotheses of \Cref{sec:input}, one has
$m_{\boldsymbol{n}}\neq2$.
\end{theorem}

\begin{proof}
\Cref{prop:d2-remaining-form} reduces the case
$m_{\boldsymbol{n}}=2$ to
\eqref{eq:d2-cyc-start}.  \Cref{prop:d2-cyclic-preparation}
gives \eqref{eq:d2-cyc-prepared} and \eqref{eq:d2-cyc-xN}, which
\Cref{prop:d2-cyclic-separation} excludes.
\end{proof}

\begin{proof}[Proof of \Cref{thm:main}]
If $(d\boldsymbol{n})_0=0$, the reduction theorem
\cite[Theorem~1.1]{PortillaReduction} gives
$m_{\boldsymbol{n}}\in\{2,3,4\}$.  \Cref{thm:dfour} excludes
four.  \Cref{thm:dthree,thm:dtwo} exclude three and
two, respectively.
\end{proof}

\printbibliography

\end{document}
